\clearpage
\thispagestyle{empty}
\pdfbookmark[0]{Modern edition: authorship and sources}{modern-edition-notice}
\begingroup
\pagestyle{empty}
\footnotesize\raggedright
\enlargethispage{2\baselineskip}
\setlength{\parindent}{0pt}\setlength{\parskip}{.5\baselineskip}
{\large Modern edition: authorship and sources\par}
Historical author: Carl Friedrich Gauss. The historical editors, correspondents, contributors and issuing institutions retain the attributions printed in the source and reproduced in the historical text. Their writing is not attributed to an AI system.

Source edition: Carl Friedrich Gauss, \emph{Werke}, Band VII, 1906, containing \emph{Theoria motus corporum coelestium in sectionibus conicis solem ambientium}, Hamburg, 1809. Digital witness: \href{https://gdz.sub.uni-goettingen.de/id/PPN236008730}{Göttinger Digitalisierungszentrum, PPN236008730}.

This original-language partial edition contains the historical titles, complete Praefatio and Contenta, Articles 1--\GaussCurrentLastArticle, the Tabulae part-title and complete Tables I--\GaussCurrentLastTable, Figures 1--\GaussCurrentFigureCount, \GaussCurrentBemerkungenScope, and the Zusätze part-title, Nachlass heading and complete items 1--\GaussCurrentLastZusatz, through [18.], \GaussCurrentZusatzRemarksScope, including Oriani’s corrections, the historical note and reply by Gauss, the ninth and tenth items’ geometric proofs and editable orbital and spherical diagrams, the eleventh item’s historical correction notice, the twelfth item’s determination of probable errors in least-squares results, the thirteenth item’s relations between initial state and orbital elements, the fourteenth item’s correction of planetary elements from four oppositions, the fifteenth item’s differential formulas with a Pallas example, and the sixteenth item’s discussion of celestial zodiacs with Goldschmidt’s Iris table. The part-title Zur parabolischen Bewegung, Nachlass und Briefwechsel heading and complete notes 1--\GaussCurrentLastParabolicNote{} on elapsed time from two radius vectors and their chord, the inverse calculation of the chord, determination of the orbital elements, the triangle-to-sector relation with its logarithm table, and Gauss’s letters to Olbers of 7 and 13 January 1815 with a worked calculation and the subsequent reduction of the cometary-orbit problem, followed by the general theory of cometary-orbit calculation and its full example for the second comet of 1813, and the letters to Olbers of 29 May and to Bessel of 24 June 1815 concerning the proposed supplement and parabolic-motion table, are also included. The Nachlass und Veröffentlichung heading, the second part on the true anomaly in parabolic motion, and explanatory Articles 1--\GaussCurrentLastTabulaNovaArticle{} of Tabula nova motus parabolici are included, completing its explanation, together with \GaussCurrentTabulaNovaTableScope. \GaussCurrentPostTabulaNovaScope{} are also included. \GaussCurrentCeresScope{} are also included. \GaussCurrentPallasScope{} are also included. The remaining contents of Volume VII are outside this edition. The second fragment [11] of Nachlass VI, printed pages 561--564, is not included: eight lower-table row labels on page 563 remain unresolved in the available witnesses. The included letter [1] of Nachlass VII begins at the page 565 position and ends on page 567 before [2]. The sequence therefore has an explicit gap between the first VI [11] and VII [1]. Historical wording, mathematical expressions and notes are retained. Modern fonts and reflow represent content rather than historical type identity. Qualified readings are described in the accompanying README.

AI-transcribed text, AI-produced source comparisons and typesetting, and AI-created editable mathematical drawings: OpenAI GPT-6 Astra (Ultra). This credit concerns the modern encoding and does not replace historical authorship. No human or independent review is claimed.

In Tabula nova, Pars II, printed page 359, the final digit of argument 5,354 and the zero in 20,228 at argument 5,294 are qualified readings from incomplete contours. The point above the final 6 of the value at 5,012 records an observed mark without assigning mathematical meaning; its printing intent is uncertain. On printed page 360, the filled superscript mark after 67 in the second block’s first row and the final digit of argument 5,620 are read as zero with support from corresponding forms and the comparison witness; their principal interiors are not fully visible. In Pars III, printed page 367, the seconds digit at argument 32200 is read as 8 from a surviving lower loop and crossing stroke; its upper loop is incomplete and the reading remains qualified. In the Encke letter, printed page 368, a lowered colon represents two vertically separated contours beside the first denominator A. Their punctuation identity and printing intent are uncertain; damaged type or printing/imaging irregularity remain alternatives. No mathematical meaning is assigned to the mark. The comma following the logarithmic interpolation fraction on printed page 369 is represented at its observed fraction-bar height. In the letter of 1 April 1843, the first calculation on printed page 372 reads the dense exponent over q as 3/2 with qualified confidence; the clearer earlier definition supports this reading, while a faint comparison witness does not resolve every stroke. A lowered period after Gleichung on printed page 373 represents a detached point at descender height; its punctuation identity and printing intent remain uncertain. In the editorial Bemerkungen on printed page 374, the comma following the sector-to-triangle fraction is represented at its observed fraction-bar height; its printing intent remains uncertain.
 In Ceres section [2], printed page 378, the first denominator exponent in the definition of R is read as 3 with support from the comparison witness; its dense principal contour also admits a reading of 2 or filled/damaged type. The final prime and closing parenthesis in the second root are qualified readings from bent and fragmentary contours. The raised formula punctuation on printed pages 378–381 represents observed positions without added mathematical meaning; printing intent remains uncertain.
 In Ceres section [3], printed page 382, the two raised commas retain their observed positions. A plus with a detached lower stroke is represented by a normal editable plus without a claim of continuous ink; the cause of the interruption is uncertain. The printed unprimed M in the Delta-w definition is retained.
 In Ceres section [4], the point above the equals sign in the first geometric-mean relation and the low point after P in the C relation retain their observed positions without added mathematical meaning. Their printing intent is uncertain. The raised comma after the limiting P expression and the printed root with following fractional exponent are retained.
 In Ceres section [5], the final M and closing parenthesis of the first integrated product are supplied by the Internet Archive witness, native page 397; the corresponding GDZ expression ends visibly at the preceding 2. The witness parenthesis survives as a lower curved stroke, so its reading is qualified. The printed coefficients 35/24 and 11/8 and the expression 3entM are retained without mathematical emendation.
 In Ceres section [6], the two terminal M letters in the second velocity row are supplied by the Internet Archive witness, native page 399; both endings are lost at the right edge of the GDZ image. The adjacent closing parentheses and final brace are qualified reconstructed readings: surviving fragments, the opening delimiters and the recovered expression support them, but the complete glyph contours are not resolved. Ordinary and Fraktur B, the unprimed e in the fourth velocity result, all 98 leader points and six underlines follow the principal without mathematical emendation.
 In Ceres section [7], the first f coefficient is 1/24, and the integration constants retain the printed primes. One horizontal mark above a primed tilded omega is represented by an additional bar without assigning it mathematical meaning. The a-like Greek alpha reading has moderate character confidence: its low right foot differs from the curled Latin a, while the recovered reading and its pairing with beta provide supporting context. Several omega-primes contact closing-parenthesis contours; modern separate glyphs do not assert separate historical ink. Raised terminal punctuation and twelve individual leader points retain their observed positions.
 In Ceres section [8], a compact detached point below and to the right of the first sine identity is represented by a lowered period; its punctuation identity and printing intent remain uncertain. The terminal point of the complete periodic-latitude formula is represented at fraction-bar height without a multiplication meaning. A small irregular rising contour below the minus in XI is not confidently a separate sign. Incidental ink, damaged type, printing or imaging irregularity remain possible.
 In the Ceres perturbation-table letters, short horizontal strokes above the left of N in Nov. and above the right of n in angezeigten are represented by editable short rules. Raised points between die and VII. and after und before C,D retain their observed positions. These readings have high contour and position confidence but uncertain punctuation identity or printing intent; no additional linguistic or mathematical meaning is assigned. Irregular contours near ist, als, und in the formula array, gehen and Methode are not confidently separate signs. Incidental ink, damaged type, printing or imaging irregularity remain possible. The Ceres, Jupiter, Pallas, opposition and ascending-node symbols retain their distinct identities; modern fonts do not reproduce historical type.
 In the letter of 10 May 1805, the final astronomical symbols are Pallas and Juno, read from their pointed-loop and radiating-star heads with lower crosses. The opening parenthesis before denn survives in separated upper and lower curved strokes; its editable form records the recognizable glyph without asserting continuous ink. A tiny partly hollow angular contour above the left of w in werde is not confidently a separate linguistic sign. Contour and position confidence is high; the cause of the parenthesis interruption and the detached mark's printing intent remain uncertain. Incidental ink, damaged type, printing or imaging irregularity remain possible.
 In section [2] of the second Ceres method, the first reciprocal relation on printed page 403 has the fractional exponent encoded as 8/2 from an 8-like contour, with moderate numeral confidence. The Internet Archive comparison, leaf 413, also appears 8-like but does not resolve every stroke. A filled or damaged 3 is an alternative supported by the preceding mathematical expression and the retained transcription; that contextual inference is not substituted for the qualified visual reading. The two parentheses around S. 381 have interrupted contours. The derivative comma and reciprocal-relation point retain their observed fraction-bar heights. Detached marks below entwickeln and below/right of the final differential are not confidently separate signs. Their printing intent and the cause of interrupted contours remain uncertain; incidental ink, damaged type, printing or imaging irregularity remain possible.
 In section [3] of the second Ceres method, the comma following the constant fraction retains its observed fraction-bar height without added mathematical meaning. A compact oblique mark before vgl in the historical note, a shallow two-winged contour above the first odd-case minus and a compact point well to the right of Formeln are not confidently separate signs. Their contours and positions are clear, but their printing intent remains uncertain; incidental ink, damaged type, printing or imaging irregularity remain possible.
 In section [4] of the second Ceres method, the beta-fraction semicolon retains its observed fraction-bar position. A compact angular contour after the logr-prime table header and a tiny branched contour below the final printed-page405 formula are not confidently separate signs; incidental ink, damaged type, printing or imaging irregularity remain possible. The last first-row expression on the right of the wide printed-page406 table is transcribed through +M without an added closing parenthesis: that character is not confidently resolved at the principal edge, and the faint corresponding edge of the retained Internet Archive witness does not settle it. This is a qualified boundary reading, not a claim of a historical printing error or absence of ink. The wide table uses a landscape page; its shared row associations and both complete braced headers remain intact.
 In section [5] of the second Ceres method, the first opening parenthesis has interrupted curved strokes; the editable glyph represents the recognizable reading without asserting continuous ink. A small flattened angular contour well to the right of the first printed-page407 numerical row is not confidently separate punctuation. Incidental ink, damaged type, printing or imaging irregularity remain possible.
 In the letter of 2 July 1805, interrupted curved contours in formula A and the four-row group in B are represented by recognizable closing parentheses and a grouping brace without asserting continuous ink. Detached angular contours above-left of the i-dot in the +3,59 sine term and below-left of the node in the -4,56 term are not confidently separate signs. Incidental ink, damaged type, printing or imaging irregularity remain possible. All eighteen underlines, eight grouping braces, sixty-four individual leader points and the printed numerical distinctions are retained without mathematical emendation.
 In Brendel’s closing remarks on Ceres, a detached angular contour above the o/r gap in the second occurrence of Correspondenz is not confidently a separate linguistic sign. Incidental ink, damaged type, printing or imaging irregularity remain possible. Brendel’s account of his reconstruction from Gauss’s numerical work is retained in full.

 In the opening Pallas correspondence, the Pallas, Ceres, opposition, conjunction, Jupiter and ascending-node symbols retain their distinct identities. Three hundred and twenty-two individual omission points and all twelve historical notes are retained. A compact detached angular contour below verschiedene and a small irregular contour beside the final s in Kunstgriffs are not confidently separate linguistic signs. Their contours and positions are clear, but their printing intent remains uncertain; incidental ink, damaged type, printing or imaging irregularity remain possible.

 In the letter of 30 November 1810, a compact irregular branched contour between vorher and integrirten is not confidently separate punctuation. Its position and contours are clear, but printing intent remains uncertain; incidental ink, damaged type, printing or imaging irregularity remain possible. The historical account of a sign error is retained without using it to alter the preceding letter’s numerical values.

In the letters of 13 December 1810 and 6 January 1811, a short oblique contour after hier, a thin rising stroke above-left of the date digit 6 and a compact irregular blob above the gap after Pallasstörungen are not confidently separate linguistic signs. Their contours and positions are clear, but printing intent remains uncertain; incidental ink, damaged type, printing or imaging irregularity remain possible.

In the letter of 12 August 1811, a compact irregular contour below the first N in Neigung is not confidently a separate linguistic or mathematical sign. Its contour and position are clear, but printing intent remains uncertain; incidental ink, damaged type, printing or imaging irregularity remain possible. Both ascending-node signs and the letter-spaced phrase in irgend einer Form are retained.

In the letter of 5 April 1812, a detached V-shaped contour above and slightly left of the ordinary i-dot in mir is not confidently an independent linguistic accent. Its contour and position are clear, but printing intent remains uncertain; incidental ink, damaged type, printing or imaging irregularity remain possible. The distinct Jupiter, Pallas and Ceres signs, the letter-spaced phrases in the letter of 5 May 1812 and the printed heading Olrers on 12 May 1812 are retained.

In the letter of 2 July 1813, a detached multicomponent vertical contour before Theil is not confidently a separate linguistic or mathematical sign. Its contour and position are clear, but printing intent remains uncertain; incidental ink, damaged type, printing or imaging irregularity remain possible. The earlier letter of 8 April retains its distinct Jupiter, Pallas, Saturn and Mars signs, decimal-seconds value and letter-spaced Probe.

In the letter of 23 April 1814, detached contours after nothwendig, above beinahe, above the lower longitude-row dash, to the right of the first mean-error value and above the line-end har- in harmonirten are not confidently separate linguistic or mathematical signs. Their contours and positions are clear, but printing intent remains uncertain; incidental ink, damaged type, printing or imaging irregularity remain possible. Both row-end dashes, the fractions 1/43 and 1/60, and all individual table leaders are retained.

In the letters of 18 May and 15 June 1814, the comma after the verbal root fraction and the point after the first mass-correction fraction retain their observed raised positions, without added mathematical meaning. In the June letter, tiny detached contours below entgegengesetzte and the line-break sinn in Eigensinn are not confidently linguistic signs. In the letter of 25 September, two tiny flattened components below neugierig are likewise not confidently punctuation. Contour and position confidence is high; printing intent remains uncertain, and incidental ink, damaged type, printing or imaging irregularity remain possible. Both numerical grouping braces and the historical note’s opposition sign are retained.

In the letter of 31 December 1814, the second table heading is read as a mass increase of 1/64. The first denominator digit is dense and partly joined to the fraction bar in the principal; its identity remains qualified. The Internet Archive comparison, leaf 435, printed page 425, has a recognizable 6 supporting that reading. An 8-like principal contour, filled or damaged type and imaging irregularity remain alternatives. The final 4 is distinct. This reading is not a recalculation of the historical values. In the letter of 8 January 1816, the pointed loop and lower cross identify the printed prize symbol as Pallas.

The letters of 24 July 1816 and 15 February 1817 retain the Sun/Jupiter mass ratio, the opposition sign and the pointed-loop Pallas sign. In the letter of 2 December 1817, the mass correction is read as (+1/46): the principal denominator is dense and partly joined to the fraction bar, while the Internet Archive comparison, leaf 439, printed page 429, shows recognizable 4 and 6. The principal digit identities remain qualified; filled or damaged type and imaging irregularity remain alternatives. This is not a recalculation. The four points after the fifth-row latitude and the high point after the ninth-row longitude are retained; the high point has qualified punctuation identity and printing intent, without added mathematical meaning. Irregular detached contours below the gap before meiner and above the repeated Breite heading are not confidently separate linguistic or mathematical signs. Incidental ink, damaged type and printing or imaging irregularity remain alternatives.

The letter to Encke of 25 March 1818 retains all three upright quotation openers, the seven-row table, 68 individual leader points, four points between 1803 and 1817, and the printed omissions of angle units. A detached shallow stroke below the final t of nicht is not confidently a separate linguistic or mathematical sign. In the letter of 17 March 1822, a small detached angular contour above the receiving word verlässigen is likewise not confidently an accent. Incidental ink, damaged type and printing or imaging irregularity remain alternatives. The letters through 4 October 1834 retain all fifty omission points and the literal brace pair around Encke’s letter.

Hansen is the author of the letter to Gauss of 7 February 1843, including its integral, recurrence relations and numerical example. The second-iteration p1 and r1 double primes have dense joined contours in the principal; the Internet Archive witness, leaf 444, printed page 434, shows two distinct strokes. The principal reading remains qualified, with filled or damaged type and imaging irregularity as alternatives. The commas after the initial integral and the limiting quotient retain their fraction-bar height without additional mathematical meaning. A hooked multicomponent contour before dT, a shallow contour above-left of der and a point above the left part of medium are not confidently separate signs. The letter of 13 October 1834 likewise has detached contours above the gap after gefallen, above ich and below Mittheilung that are not confidently linguistic punctuation. Contours and positions are clear, but printing intent remains uncertain; incidental ink, damaged type and printing or imaging irregularity remain possible. All printed numerical values, both differing final results, the six table rules and the literal braces are retained without recalculation.

In the letter of 11 March 1843, a diagonal extension touches the final n of Von, and a short flattened contour lies above-right of the period after machen. Neither is confidently separate punctuation. The opening parenthesis before der has interrupted curved strokes; its editable form records the recognizable glyph without claiming continuous ink. Contours and positions are clear, but causes and printing intent remain uncertain; damaged type, incidental ink and printing or imaging irregularities remain possible.

In the letter of 21 March 1843, the prose fractions read 48/47 and 43/42; the three numerical comparison rows and the five points in the signed day interval are retained. Gauss’s subsequent statement about the Saturn perturbations remains in its historical position. Two detached contours below and to the right of the closing section rule are not confidently separate ornaments or signs. Their positions are clear, but printing intent remains uncertain; incidental ink, damaged type and printing or imaging irregularities remain possible.

The opening of the Pallas Exposition preserves its English epigraphs and the printed Thomson and Boyle attributions. In the Thomson citation, the second digit of 1346 has a dense principal contour; the reading 3 is supported by the recognizable digit in the Internet Archive witness, leaf 449. Filled or damaged type remains an alternative for the principal contour. The unaccented a after rapporté and the spelling connait are retained. Detached contours below-left of Boyle, above the a in connait, below the receiving naison and below réel are not confidently separate signs; incidental ink, damaged type and printing or imaging irregularities remain possible.

In Article 1 of the Pallas Exposition, the printed xdy term in the squared expression and the u/r cubed term are retained. The exponent after du in the first FF-chain member is read as 2 with moderate principal glyph confidence: its dense contour also permits damaged or filled type and a 3-like reading. The retained transcription and the quadratic expression support 2, but the mathematical inference is not direct ink evidence. Three terminal points and two commas retain their observed fraction-bar heights without added meaning. A low stroke beneath the opening r and detached spots above the AA-chain and far left of equation 14 are not confidently separate signs. Their contours and positions are clear; incidental ink, damaged type and printing or imaging irregularities remain possible.

In Article 2 of the Pallas Exposition, three commas, two semicolons and two terminal points retain their fraction-bar heights without added meaning. A low stepped contour between intégrer and cette, a compact diagonal spot far left of the E in equation 19, and a shallow multicomponent contour between on and aura are not confidently separate signs. Their contours and positions are clear, but printing intent remains uncertain; incidental ink, damaged type and printing or imaging irregularities remain possible. The two mass labels retain four individual points each, and all bracketed digits and the printed grouping of the logarithm are preserved without recalculation.

In Article 3 of the Pallas Exposition, three terminal points and one comma retain their observed fraction-bar heights without added meaning. A small detached multicomponent angular contour far right and slightly below the intermediate quotient is not confidently a separate sign. Its contours and position are clear, but printing intent remains uncertain; incidental ink, damaged type and printing or imaging irregularities remain possible.

In Article 4 of the Pallas Exposition, the printed supposeront and the later supposerons are retained. A compact diagonal contour below-left of soleil, a shallow angular contour above-right of the acute accent in négatif, and a vertical multicomponent contour below-right of the first plus in AA+BB+CC are not confidently separate signs. The acute accent itself is distinct. Contours and positions are clear, but printing intent remains uncertain; incidental ink, damaged type and printing or imaging irregularities remain possible.

In Article 5 of the Pallas Exposition, both fractional exponents in the combination formula read 3/2; the upper digits have dense two-lobed contours. A small detached angular, partly hollow contour below the first grands in the pole-substitution paragraph is not confidently a separate sign. Its contours and position are clear, but printing intent remains uncertain; incidental ink, damaged type and printing or imaging irregularities remain possible.

In Article 6 of the Pallas Exposition, two terminal points retain their observed fraction-bar heights without added meaning. The closing parenthesis after 1-e in the first radius-sum numerator is a qualified reading from a small upright fragment and the surrounding expressions; the faint comparison witness does not resolve its missing curved strokes. Other parenthesis contours survive as separated curved fragments. A detached angular point far right of the final introductory sentence is not confidently a separate sign. Contours and positions are clear, but printing intent and the causes of interruptions remain uncertain; incidental ink, damaged type and printing or imaging irregularities remain possible.

In Articles 7–8 of the Pallas Exposition, nine commas and one terminal point after differential quotients retain their observed fraction-bar positions without added meaning. In Article 7, tiny detached angular marks above-right of Mais and below instant are not confidently separate signs. Contours and positions are clear, but printing intent remains uncertain; incidental ink, damaged type and printing or imaging irregularities remain possible. The spaced phrase ellipse osculatrice retains its emphasis in modern editable type.

In Article 9 of the Pallas Exposition, two commas and one terminal point after the three force fractions retain their fraction-bar positions without added meaning. The minus before zeta cos ZL has an attached descending irregularity; the closing parenthesis after Bx in the final numerator survives as separated curved fragments. Editable minus and parenthesis represent recognizable glyphs without asserting continuous ink. Contours and positions are clear, while causes and printing intent remain uncertain; incidental ink, damaged type and printing or imaging irregularities remain possible.

In Article 10 of the Pallas Exposition, the opening parenthesis before the reference 4, 5, 6 survives as separated upper and lower curved fragments. The energy differential’s terminal point retains its fraction-bar position without added mathematical meaning. The recognizable parenthesis does not assert continuous ink. The exponent glyphs in the fourth-power term of equation 64 and the cubed terms in equation 65 and the following phi differential have dense ink; their observed forms are retained. Incidental ink, damaged type and printing or imaging irregularities remain possible explanations of irregular contours.

In Article 11 of the Pallas Exposition, the opening parenthesis before 1 in the first substituted dp/r relation has separated curved strokes. Its editable form represents the recognizable glyph without asserting continuous ink. A tiny detached oblique contour below the left side of equation 68 is not confidently a separate mathematical sign. Contours and positions are clear, but printing intent and the cause of interruptions remain uncertain; incidental ink, damaged type and printing or imaging irregularities remain possible.

In Article 12 of the Pallas Exposition, three terminal points and two commas are represented at fraction-bar height without added mathematical meaning. Several parentheses and the closing brace in the five-member chain have interrupted curved strokes. The final cosine argument has only partly surviving parentheses: the opening curve remains a qualified contextual reading, and the faint comparison witness does not resolve the missing contours. A tiny detached contour below the first long dL expression is not confidently a separate mathematical sign. Observed positions and contours are clear, but printing intent and the cause of interruptions remain uncertain; incidental ink, damaged type and printing or imaging irregularities remain possible.

In Article 13 of the Pallas Exposition, a short detached, slightly stepped horizontal stroke follows the lower right foot of r in pour. A lowered hyphen represents its observed direction and low position; its identity and printing intent remain uncertain, without an assigned linguistic or mathematical meaning. A small oblique contour above the final integral sign is not confidently a separate mathematical sign. Incidental ink, damaged or extended type and printing or imaging irregularities remain possible.

In Article 15 of the Pallas Exposition, the terminal point after the third force quotient is represented at fraction-bar height without added multiplication meaning. A small partly hollow irregular contour in the whitespace below the final distance formula is not confidently a separate linguistic or mathematical sign. Observed contours and positions are clear, but punctuation identity and printing intent remain qualified; incidental ink, damaged type and printing or imaging irregularities remain possible.

In Article 16 of the Pallas Exposition, exponent 2 in the final denominator of R/mu is a qualified reading of the principal contour. Its lower body is densely filled; the faint comparison witness does not resolve the digit. Exponent 3 remains a contextual alternative. The terminal point after the W quotient is represented at fraction-bar height without added multiplication meaning. Observed contours and positions are clear, while glyph identity and printing intent remain qualified; incidental ink, filled or damaged type and printing or imaging irregularities remain possible.

In Article 17 of the Pallas Exposition, a short detached contour above the gap between d and u in du moins is not confidently a separate linguistic sign. Its position and contour are clear, while printing intent remains uncertain; incidental ink, damaged type and printing or imaging irregularities remain possible.

In Article 18 of the Pallas Exposition, the inner closing parenthesis in the third-derivative pair of series III has interrupted curved strokes; the editable parenthesis represents the recognized glyph without asserting continuous ink. Three terminal quotient points are represented at fraction-bar height without multiplication meaning. Small detached contours below fonction, beyond the third-difference definition and before exprimer are not confidently separate linguistic or mathematical signs. Their contours and positions are clear, while printing intent remains uncertain; incidental ink, damaged type and printing or imaging irregularities remain possible.

In Article 20 of the Pallas Exposition, a low point follows supposer in the principal and the faint comparison scan; its representation as a period does not assert grammatical intent. The accent above omega has a thick, shallow contour with an upturned right end. Its representation by a tilde is a qualified glyph reading; a bar-like interpretation remains possible, and the faint comparison scan does not settle the printing form. A small extension and speck below précision and a detached point below-left of the closing rule are not confidently additional linguistic or mathematical signs. Their visible positions do not establish whether they arose from type, ink, printing or imaging.

In Article [21] of the first draft of the fourth Pallas Exposition section, a compact detached contour above-right of du in longitude du noeud is not confidently a separate linguistic or mathematical sign. Its observed position and irregular boundary do not establish its printing intent; incidental ink, damaged type and printing or imaging irregularity remain possible. The faint comparison witness, leaf 475, printed page 465, does not settle that intent.

In Article [22], five L-prime and closing-parenthesis pairs have touching or merged upper contours. Both glyphs are retained as a qualified reading supported by the explicit primed longitude and the surrounding notation; separated historical ink is not asserted. The opening parenthesis after dt has interrupted curved strokes. Two compact detached contours above-right and below-right of est near the end are not confidently additional linguistic or mathematical signs. Incidental ink, damaged type and printing or imaging irregularities remain possible. The faint comparison witness, leaf 476, printed page 466, does not independently resolve all these contours. The two raised oval zeros after 360 and the baseline point after dt are retained.

In Article [23], the potential denominator exponent is represented as 8 from its dense two-lobed principal contour. This is a qualified glyph reading: an ink-filled or damaged 3 remains possible. The retained earlier transcription reads 3, and the faint comparison scan has a more open 3-like contour; the mathematical context also favors 3, without settling the principal printing form. Three shallow omega accents are represented by tildes, with a bar-like alternative retained. Two L-prime and closing-parenthesis pairs, and the bars adjoining parentheses in the mean-longitude formula, have touching contours; the represented characters do not assert separated ink. The raised hyphen-shaped stroke between au and tems and the incomplete descending comma after seconde fois retain their positions without asserting printing or grammatical intent. A detached contour above-left of loin is not confidently a separate linguistic or mathematical sign. Incidental ink, damaged type and printing or imaging variation remain possible. The four roots are unbarred. Raised commas and the raised terminal point in the force equations are punctuation; the terminal point is not assigned multiplication meaning.

In Article [24], the closing parenthesis in the historical note call has interrupted upper and lower curved strokes. The recognizable character is retained without asserting continuous ink. The note opens with a bracket, asterisk and parenthesis, and closes with a bracket after its terminal point. In the page reference 557–559, the final 9 is a qualified reading of a larger upper loop and dense lower returning contour, supported by the retained earlier transcription; an 8-like appearance remains possible. The extremely faint comparison note does not resolve the digit. Damaged or ink-filled type and printing or imaging variation remain possible.

In Article [25], a small detached stroke above the gap between valeurs and justes is not confidently identifiable as a separate linguistic sign. The faint comparison does not establish its origin; incidental ink, damaged type and printing or imaging variation remain possible. The seven superscripts of Q are read as zero, with several dense or joined contours in the principal. The first draft ends with the printed unfinished sentence, five omission points and a closing rule.

In Article 21 of the second draft, the period 360 retains its raised oval zero. Small detached marks above or just before the opening On aura a on printed page 470 are not confidently identifiable as linguistic accents or a mathematical prime; the variable is read as unprimed a. The faint comparison does not determine their origin. Isolated marks outside the displayed expressions have no established operator meaning. Small gaps in the central coefficient fraction bars are retained as recognizable bars without asserting continuous ink. Incidental ink, damaged type and printing or imaging variation remain possible.

In Article 22 of the second draft, two small dense marks to the upper and lower left of U on printed page 471 are provisionally read as upright e. Their existence and positions are clear; their glyph identity is qualified. The comparison preserves the upper e-like contour but loses much of the lower one. No limiting-value or other mathematical operation is assigned to these marks. Damaged type and printing or imaging variation remain possible. The isolated point well to the right of the T label and the small mark below and to the right of the closing rule have no established operator, label-punctuation or ornament meaning; the faint comparison does not resolve their origin. Baseline product points, repeated ee, the literal placement of powers after the parenthesized trigonometric half-arguments, tang and cotang, and both raised oval zeros after 90 are retained.

In fragment [1] of the first calculation of the special perturbations of Pallas by Jupiter, the rho-denominator exponents in formula groups 1 and 2 are read as 3 with qualification. Both principal contours are densely filled; the comparison witness supplies a clearer 3-like form in group 1 and a faint 3-like fragment in group 2. This reading is also consistent with the mathematical context, but is not an unambiguous direct reading of the principal. The retained reading 2 and the filled, 8-like appearance remain alternatives. A detached point above v in the sine factor of the omega differential is retained in position, without assigning a derivative or other mathematical meaning to it. The faint comparison does not resolve its origin. Damaged type and printing or imaging variation remain possible. Latin w-prime, the ascending-node sign, the shallow omega accents, the unbarred root, baseline product points, repeated-letter products, and the literal power position in cos phi-squared are preserved. Modern fonts and reflow represent the content rather than historical type identity.

In fragment [2] of the first calculation, Aquinoctium in the opening and Äquinoctium in the later reduction paragraph retain their distinct printed forms. The note has only n visible between Gauss and Olbers in both witnesses; the possible missing a is not supplied. The last latitude is read as 19,8 with qualification because the upper loop of 8 is interrupted. The densely filled raised signs after 121 and 14 are read as zero with qualification. Three detached points above table leaders and a point above-left of the first System I label are retained in position, without assigning mathematical or punctuation meaning. The comparison confirms the last-row leader point but does not resolve all the other points. The detached contour above aus and other isolated marks do not confidently establish additional linguistic signs. Damaged type and printing or imaging variation remain possible. Omitted unit marks, the ascending-node signs, and the literal table entries are retained. Modern fonts and reflow represent content rather than historical type identity.

Fragment [3] of the first calculation retains the staggered first sums, the double sums on the main rows, and the extra terminal double sum in its own column. The literal header 1830+, the terminal sum -19,844, and the entry +411 are retained without numerical emendation. The principal supports -285,682 and +64,971 in the cumulative columns; the comparison image is too faint to resolve the latter independently. The inclination ends in 4, and the interrupted digit in -9,144 is supported by the comparison. The first opposition epoch retains a raised d followed by 0 without a supplied decimal point. A detached raised stroke after 25 46 42, a low point following 55,737, and two marks near the comma of 6,362 are represented in position without assigning an operation, unit or punctuation role. The broken contour beside the minus in -2,620 does not confidently establish another sign. Other isolated ink and interrupted contours likewise remain uncertain; damaged type and printing or imaging variation are possible. No closing square bracket is supplied before fragment [4]. Modern fonts and reflow represent content rather than historical type identity.

Fragment [4] retains the separate words auseinander gesetzten, the tangent abbreviation tang, the literal p and r labels, the shared numerical separators, all omitted units, and the closing wording Dass gefundene. Its first observation epoch has a comma beneath the raised d, followed by 019120. The fourth mean anomaly ends in 57,16 and the second corrected epoch in 36,35; the dense final contours do not warrant replacing either 6 by 8. Raised oval units are represented by superscript zero, including filled principal contours supported by open forms in the comparison image. A short stroke after Roman V, a tiny mark above the fourth observed latitude plus, a point below the comma in \(768'',2969\), and a point below c are represented in position without assigning an operation, unit or punctuation meaning. The faint comparison does not independently resolve all these marks. Interrupted parentheses and isolated ink do not confidently establish further signs; damaged type and printing or imaging variation remain possible. Modern fonts and reflow represent content rather than historical type identity.

Nachlass IV, fragment [1], retains the literal 829 in the true-anomaly column and the space without a comma in 0 44746. The faint comparison permits a broken 3 or an 8-like contour for the first digit of 829 and does not establish a silent correction to 329. The fourth reference longitude retains the qualified final-nine reading in 226 38 39,24; its damaged contour also permits a three-like appearance. The first longitude retains a double prime despite its dense second contour. Raised oval units are represented by superscript zero, including the filled first true-anomaly unit supported by an open comparison form. A short stroke above the comma of 0,32129, a curved stroke above the minus of the sum −30,8747, and a short stroke after the last digit of the sum −151,831 are represented in position without assigning another operation or punctuation meaning. Faint comparison printing does not independently resolve these marks; incidental ink, damaged type, printing and imaging variation remain possible. The headers preserve 50dn and d50n in their respective panels, with all staggered sums, initial zeros, supplied and omitted units, and the separate terminal double sum. The continuation on printed page 487 belongs to this fragment and precedes its closing observation in the present reflow. Modern fonts and reflow represent content rather than historical type identity.

Nachlass IV, fragment [2], retains 129 59 40,31 in the fifth omega column and 0,415 2758 in the second p column; the principal and comparison support these readings. The retained earlier transcription has final 6 in the latter value. The first epoch has a comma beneath its raised d, followed by 02; the faint comparison does not independently resolve that comma. Raised oval units are represented by superscript zero, and interrupted parentheses remain recognizable parentheses without asserting continuous ink. A short stroke beneath the node in the fifth equation and a small point immediately before von are represented in position without assigning a mathematical operation or punctuation intent. The faint comparison does not resolve the stroke and does not distinctly reproduce the point. An isolated mark above-left of the final residual list is not confidently a separate textual sign. Incidental ink, damaged type and printing or imaging variation remain possible. All seven numerical columns, omitted units, shared separators, seven equations and terminal residuals retain their historical content. Modern fonts and reflow represent content rather than historical type identity.

Nachlass V, fragment [1], retains the historical spelling Entwickelung, tang, log hyp, the astronomical node symbols, eleven ditto signs and the repeated-letter products. Three heavy denominator exponents in the nine-definition group are read as 3; the faint comparison supports the T-denominator reading more distinctly than the two others, and filled 8-like contours remain a visual alternative. The terminal point after the alpha fraction is represented at fraction-bar height without multiplication meaning. An upper point and a lower short stroke immediately before Pallasbahn are represented in position without assigned linguistic or mathematical meaning; the comparison does not distinctly reproduce them. An isolated point far right of also is not confidently a textual sign. Shallow omega accents and interrupted parentheses are represented by their recognizable glyphs. Incidental ink, damaged type and printing or imaging variation remain possible. Modern fonts and reflow represent content rather than historical type identity.

Nachlass V, fragment [2], retains Werthe and Werte in their respective passages, nemlich, roman J, all half-endpoint factors, fourteen numbered groups and twenty-five baseline leader points. Group 4 has m alpha without an n subscript; the later even-case cosine identity retains mu/2. Neither is silently emended. Three detached points, near der, beneath sin and upper-right of the final L in a definition, are represented in position without assigned linguistic, operational or prime meaning. The faint comparison does not distinctly resolve their printing intent. A baseline point after na before etc. is retained without multiplication meaning. Incidental ink, damaged type and printing or imaging variation remain possible. Modern fonts and landscape reflow represent content rather than historical type identity.

Nachlass V, fragment [3], retains the printed specialisirt, nemlich, roman J, every half and quarter factor, the raised zero in roman II and all unbarred roots. The last coefficient in numbered group 3, beta22, reads minus one-half J6. Four terminal numerators in roman IV read E2 minus E10 minus twice E6. Detached paired marks, low and vertical strokes, the upper-right stroke beside D0 and the upper-right mark beside J3 are represented positionally without assigned prime, linguistic or operational meaning. Raised terminal points and the low point beside E3 have no multiplication meaning. Incidental ink, damaged type and printing or imaging variation remain possible; modern fonts and reflow represent content rather than historical type identity.

Nachlass V, fragment [4], retains the printed numerical values, including the two central one-half factors on page 502, the value 1,68124 in the c/48 line, the order of the B-prime angle on page 504, and the negative zero and final cosine term in the complete series. Unbarred roots, primes, raised zeros, dividing rules, starred references and the observed baseline leader points are retained. The closely spaced third and fourth points in the final c-prime/48 leader have no assigned operational meaning. Qualified isolated contours may reflect ink, damaged type, printing or imaging; modern fonts, spacing and reflow do not reproduce historical type identity.

Nachlass V, fragment [5], retains two detached upper/lower pairs: one before the epsilon row’s plus in group 1 and one between the plus and q12 in group 3. Each pair is represented by two editable components without assigning colon, division or other operational meaning. The last row of group 4 has dense, filled indices read as nu with qualification; the faint comparison does not resolve their original contours. The separated index digits in the coefficient-list r1 0 and in q2 1 of the delta21 relation retain their gaps without supplied commas. A faint or lost comma remains possible. Isolated marks above darum, near mögen and the coefficient examples, beside gamma10 and below zeta21 are not confidently additional linguistic or mathematical signs. Incidental ink, damaged type and printing or imaging variation remain possible. The three baseline leaders retain 47, 32 and 55 individual points. The initial A00 in the sine example, the undoubled p10 and the absence of a plus after the first constant example’s ellipsis are retained. Modern fonts, component shapes, spacing and reflow represent content and position rather than historical type identity.

Nachlass V, fragment [6], retains all three value tables and three coefficient tables, their empty cells, negative zeros, implicit units and the original order of the argument differences. Two commas, in V at 240 degrees and M=210 degrees and in W at 30 degrees and M=120 degrees, are represented below the ordinary baseline. T at 337 degrees 30 minutes and M=330 degrees reads +1,583; V at 352 degrees 30 minutes and M=150 degrees reads -35,255; the sine-2 coefficient of V at M=60 degrees reads -6,280. The last digit of V at 300 degrees and M=330 degrees is read as zero in -9,910, but its open oval does not exclude a damaged 6. W at 135 degrees and M=90 degrees is read as -98,905; the principal has a 3-compatible first contour, while the reprint supports 9. The T cosine-4 coefficient at M=30 degrees is read as +2,039 and the V sine-3 coefficient at M=60 degrees as -7,766; the principal has respectively 3-compatible and closed 8-compatible terminal contours, while the comparison witnesses support the selected digits. The T sine-10 coefficient at M=120 degrees is read as +0,016 from the reprint; the principal has a closed 8-compatible contour and the faint alternative does not resolve 6 versus 8. The alternative +0,018 remains possible as printing intent. The W coefficient readings +10,983, +2,860, +8,783, +8,376, +8,808 and +13,458 retain qualified 8-shaped contours, some broken or S-like. The T cosine-8 coefficient at M=240 degrees retains the visible minus 3, a gap, and 524 without a supplied comma. Dense nu-shaped prose indices are qualified readings. Isolated contours beside the opening elements and within nominally empty table cells are not confidently additional linguistic signs, operators or numbers. Incidental ink, damaged type and printing or imaging variation remain possible. The 164 opening leader points and final closing bracket are retained. Modern fonts, positional components and reflow represent content rather than historical type identity.

Nachlass V, fragment [7], retains the two forty-row expansions of W, their differing arguments, all seventy-eight final periodic terms and the four secular values. The cosine expansion reads -15,476 cos 2M -8,473 sin 2M with unprimed M; the corresponding sine expansion has -1,141 cos 2M prime -14,972 sin 2M prime. The logarithmic example reads log P sin P with an uppercase first P. In the first final node term, the phase retains a raised zero after 123, followed by 56 and 7 seconds, without a supplied minute sign. The closing parenthesis after the argument 2M prime minus 3M consists of three broken aligned contours in the principal; the reprint supports a closing parenthesis. A low point below the gap after M prime in the -2,052 sine term is not confidently an additional linguistic sign or operator. Detached low ink beside the node heading, a far-right contour after the phrase introducing the auxiliary angles, and a short stroke below omega in the logarithmic table likewise remain unassigned. Incidental ink, damaged type and printing or imaging variation remain possible. Filled raised-zero contours remain qualified readings. The three unbarred roots, all 307 leader points, five ditto signs, original implicit units and complete closing bracket are retained. Modern fonts, positional components and reflow represent content rather than historical type identity.

Nachlass V, fragment [8], retains the complete eighty-four-term perihelion series and both pairs of secular values. Its logarithmic table contains 71 individual leader points and a comma-shaped contour after fourteen points in the second run, followed by two further points. The comma represents the observed contour and position without an assigned mathematical or linguistic meaning; the principal and reprint both display it. The phase fields retain their sparse unit marks, including the final seconds mark after zero in the 3M prime plus M row and the fully marked angle in the 11M prime minus 4M row. The first daily secular value is -0,01310 without a supplied seconds mark. The coefficient 6,69, the coefficient -23,34 and the minus before 1,90 remain as printed. Filled raised zeros remain qualified readings. Detached marginal and low contours beside the prose, annual expression and series are not confidently additional operators, accents or punctuation. Incidental ink, damaged type and printing or imaging variation remain possible. Modern fonts and reflow represent content rather than historical type identity.

Nachlass V, fragment [9], retains the complete seventy-term series and the four secular values. In the term with coefficient -15, the phase ending after 43 is represented by one prime followed by two full-height closing parentheses. Both the principal and the reprint support these contours; the contextual reading as seconds followed by one parenthesis is not substituted. The coefficient +37 in the sixth-harmonic term is retained. The positive secular values +117,48 and +0,32163 have no supplied seconds marks. The logarithmic leaders contain eighty individual points; the slightly descending seventh point of the first run is represented as a point, with incidental ink or printing and imaging variation remaining possible. The differential expression retains its negative sign, numerator 10000000, denominator 206265, factor 2 lambda and baseline product point; the prose divisor is 206265. Detached marginal and low contours are not confidently additional punctuation or operators. Sparse angle units, qualified raised zeros and the original position of the exponent in the closing eccentricity relation are retained. Modern fonts and reflow represent content rather than historical type identity.

Nachlass V, fragment [10], retains the complete fifty-nine-term series. Both small factors in the differential equation and the logarithmic row are represented as 2/3. Their heavy denominator contours are compatible with an 8-like filled form in the principal and reprint; the faint alternate scan is compatible with 3 but does not settle the printing intent. The retained transcription and mathematical context support 3 as a qualified reading, not as proof of the printed glyph. The intermediate logarithmic value 2,23410 has neither an n subscript nor an asterisk; the reference to an asterisk in the following fragment does not supply either mark here. Only 8,74966 and 3,38137 carry n subscripts. The low irregular contour between the coefficient 17 and the following cosine is not confidently an additional sign; ink, type, printing and imaging variation remain possible. The first negative series term retains the argument M prime minus 4M without an extra sign. The opening products retain repeated aa and rr, the exponent on varphi, the shallow omega accent and the baseline product point. The seventy-seven individual logarithmic leader points, sparse angle units, printed zero fields, closing period and bracket are retained. Modern fonts and reflow represent content rather than historical type identity.

Nachlass V, fragment [11], retains all fifty-nine negative periodic terms. The opening longitude sign is represented by an editable pointed loop above a cross, with its vertical stroke entering the loop, without assigning a standardized character identity. Its contour is supported by the principal and reprint; the historical type design and printing intent remain qualified. The example retains -3051 with a seconds mark and decimal comma before 00, whereas the corresponding series coefficient is -3121,14. Its closing parenthesis has separated strokes in the principal and a complete contour in the reprint; the lower stroke is not represented as an extra comma. The first logarithmic row has an asterisk and an n subscript on 2,23410; this does not supply either mark to the preceding fragment. The two n subscripts, fifteen individual leader points, logarithmic rule, zero fields, sparse angle units and closing bracket are retained. Modern fonts, reflow and the editable sign represent observed content rather than historical type identity.

Nachlass V, fragment [12], retains both unreduced equations, six node signs and four shallow omega accents, represented by the established editable node sign and tilde-omega. Exact historical type identity and printing intent remain qualified. The thickened left end of the minus in the argument of the +2,26 term does not establish a separate operation; the actual baseline product point in the opening expression is retained. The paired enlarged readings support coefficient -166,00 and phase ending 162 54 33. All sixty-eight periodic terms, forty-two individual logarithmic leader points, six logarithmic decimal values, four n subscripts and the rule are retained. The annual secular value has a seconds mark before its decimal comma; the daily -0,04183 has no seconds mark. The final explanation continues the series paragraph. Modern fonts and reflow represent content rather than historical type identity.

Nachlass V, fragment [13], retains all seven opposition rows and their abbreviated numerical entries without expanding omitted leading fields. The anomaly rows have different unit marking: the second row retains bare spaced number fields. The article reference [10], both perturbation sequences, four logarithmic series and the stated means remain as printed. The node sign and shallow omega accent use the established editable forms; exact historical type identity remains qualified. A detached diagonal contour near the fifth final phi entry, 20 seconds comma 8, has no established mathematical or linguistic meaning and is not encoded as an additional sign; its decimal comma is separately visible. The three final quantities retain their seconds marks before decimal commas, and the period precedes the closing bracket. The terminal rule is editable. Modern fonts and reflow represent content rather than historical type identity.

Nachlass VI, fragment [1], retains the 1810 element system, the source Jupiter mass, sparse angle units and all twelve table halves. The complete-width reprint supplies the rightmost columns clipped in the principal scan, including the recorded incomplete portions of the penultimate column of the first W coefficient table. Values, explicit signs and unsigned zeros retain the source forms without numerical reconstruction. Sparse comma strings in the coefficient tables, including +1,204 in the second T table, carry no added numerical interpretation. The first T value table has the qualified ten readings in its final column supplied by the reprint. In the second T value table, the comma at row 157 degrees 30 minutes, M=285 degrees, is raised. The first V value table retains the raised terminal point at row 150 degrees, M=60 degrees; this point has no assigned multiplication meaning. Its row 337 degrees 30 minutes, M=90 degrees, retains the qualified principal 8 in +34,827 despite the reprint's 5-like contour. The second V table retains the damaged 8 reading in -19,188 at row 135 degrees, M=195 degrees. In the first W value table, +12,429 at row 292 degrees 30 minutes, M=30 degrees, retains its qualified final 9 and single comma. The second W value table retains lowered digit groups 6, 9 and 11 in the M=255 degrees column at rows 307 degrees 30 minutes, 315 degrees and 322 degrees 30 minutes; they are not converted into newly calculated values. Its -9,005 at row 135 degrees, M=270 degrees, and -10,828 at row 307 degrees 30 minutes, M=285 degrees, retain qualified principal readings against different reprint contours.

In the first V coefficient table, -6,730 at cosine 3, M=0 degrees, remains a qualified reading of a damaged digit compared with the reprint. The cluster after +0,094 at sine 5, M=345 degrees, in the second V coefficient table has no established separate sign identity. In the first W coefficient table, the cluster near +6,750 at sine 1, M=15 degrees, and the broken minus at sine 6, M=90 degrees, occur in both witnesses. Neither is numerically reconstructed. In the second W coefficient table, the constant +36,094 retains the complete plus sign; cosine 3, M=195 degrees, retains the principal's compact point in -28.847, while the reprint shows a comma-like tail. No multiplication or decimal intention is assigned to that point. The qualified final 3 readings in +23 at cosine 14, M=300 degrees, and -53 at sine 21, M=240 degrees, remain. Other isolated ink contours do not establish additional operators or punctuation.

The six closing equations preserve all factors, both occurrences of the node term where printed, the complete tang spelling, eight low inter-factor points, repeated bb and rr, and the power after phi in cos phi squared. The short omega accent retains the qualified tilde-omega reading. The broken lower stroke of the final parenthesis in the second equation is represented as a closing parenthesis. Two detached principal contours near the fourth equation have no established separate mathematical or linguistic identity. The final point precedes the tall editorial bracket that closes the opening on printed page 529. Historical signatures VII./67, 67 with asterisk, VII./68 and 68 with raised asterisk remain. Modern fonts, local positioning and reflow represent content and observed positions rather than identical historical type or page layout.

Nachlass VI, fragment [2], retains the two calculation grids, their three rules and 185 explicit leader points. The node example has no terminal point; the inclination example has one. Jupiter and Pallas longitude signs and short omega accents use editable modern font forms without claiming historical type identity. Both series retain 104 sine terms in left-column then right-column order, with the source signs and sparse angle units. Eight column starts carry full angular units; the other terms do not acquire repeated units. The coefficient 8,00 remains literal, and no additional seconds mark follows 2,47. Both node secular values and the daily inclination value carry seconds above their commas; the annual inclination value +4,9604 does not. The closing point and editorial bracket are retained. Detached contours at six series loci have no established additional sign identity: the high point after 18 in the angle 351 34 18; ink below the minus in the 5-Jupiter-minus-6-Pallas term; the contour after the 6-Jupiter-plus-49 23 13 parenthesis; the offset point above sin in the inclination coefficient 2,43; fragments near the opening parentheses of the 0,19 term with angle 214 48 17 and its following term; and the low contour below the minus in the 9-Jupiter-minus-4-Pallas inclination term. Detached marks above the inclination calculation heading, after the reprint node coefficient and beneath the reprint transformation Q are likewise not encoded as extra signs. Ink, type damage and imaging remain alternatives. Modern fonts and reflow represent content and observed positions rather than historical type or page identity.

Nachlass VI, fragment [3], retains the original cos phi squared position, the full p/r factor, the low inter-factor point before T, both node terms and the qualified omega accents. Its integration calculation has six decimal values, 55 explicit leader points and one horizontal rule. The example coefficient +653,01 carries seconds above its comma and retains its terminal point. The 160-term series follows each source left column before its right column. Four column starts carry complete angle units; the remaining terms retain sparse unit marking. The coefficient 1,00 remains literal. The daily secular value has seconds above its comma; the annual value -10,826 has none. Both historical signatures VII./69 and 69 with asterisk remain. In the 13,33 term with angle 38 29 5, the Pallas contour has a short or indistinct lower stem; Pallas remains a qualified glyph reading without a claim of identical historical type. The reprint point after 91 in the 1,13 term, raised trace before 14 in the 1,32 term and angled contour above 1,08 do not establish extra punctuation or units. The principal stroke below-left of the 0,44 leading plus and point above the minus in the 0,63 term likewise have no established additional operator meaning. Broken parentheses, the damaged Jupiter contour in the 72,38 term and the small trace beneath sin in the 0,16 term retain their incidental-contour qualifications. The reprint supports the broken principal 2 in the angle ending 20 12 20. Ink, damaged type, printing and imaging remain alternatives. Modern fonts and reflow represent content rather than historical type or page identity.

Nachlass VI, fragment [4], retains the observed u-shaped denominator and its low point before dt. The grid n has a different contour; the intended historical type, including a possible inverted or damaged n, remains uncertain, and no contextual replacement is made. The integration calculation retains six decimal values, 55 explicit leader points, one horizontal rule and the subscript n. The example coefficient +213,24 has seconds above its comma and retains its terminal point. All 123 sine terms follow the left source column before the right on each leaf. The first four longitude contours are Pallas. Four column starts have full angular units; the other 119 terms retain sparse units. The initial delta phi has no added plus. The daily secular value -0,029556 has seconds above its comma; the annual value -10,795 has none and retains its final point and editorial bracket. The low mark under 37 in the 0,22 term, detached contour beside the plus in the 0,27 term and low stroke between 30 and 5 in the final 0,11 term on printed page 548 do not establish added signs or units. On printed page 549, the low stroke after the 3,14 closing parenthesis, low contour below the 0,70 closing parenthesis, raised mark between sin and the parenthesis in the 0,11 term, and short stroke between 131 and 28 in the 1,64 term remain incidental-contour qualifications. Reprint marks beneath the 2 of 0,92 and beneath the first 2 of 202 on printed page 548 are not added symbols. The isolated principal blob below the result introduction has no established linguistic or mathematical function; its compared reprint region is blank. Interrupted parentheses, fragmented Jupiter tails and irregular digits retain their contour qualifications. Ink, type, printing and imaging remain alternatives. Modern fonts and reflow represent content rather than historical type or page identity.

Nachlass VI, fragment [5], retains the observed two-lobed, 8-shaped denominator in the opening fraction with numerator 2. The higher-resolution reprint detail shows two counters. The retained leaf transcription reads 3; a printed 3 with filled or damaged contours remains an alternative, and the intended historical type is unresolved. The factors lambda and 10000000/206265 remain unchanged. The integration grid preserves eight decimal values, 74 explicit leader points, both horizontal rules, two subscript n markers and the asterisk without an added multiplication meaning. The phrase preceding the cosine example has no added colon. The cosine coefficient is -2422; the closing sine coefficient is +2422, with the actual Jupiter and Pallas longitude signs, all angular units, final point and editorial bracket. Isolated principal ink above the p-region of Beispiel and detached reprint ink near vollständige do not establish additional accents or punctuation. Modern fonts and reflow represent content rather than historical type or page identity; the denominator reading remains qualified.

Nachlass VI, fragment [6], preserves the two paragraphs, integration calculation and all 145 perturbation-table rows in original column-down order. Each of the six original panels retains its headings; full angular units and the seconds mark above the decimal comma occur in its first row, with sparse units preserved in the remaining rows. The right panel on printed page 551, row 13, reads -0,19 in both compared scans; the retained leaf transcription has +0,19. Row 22 ends its angular triple with 5. The initial relation has no terminal point, and the final cosine example remains within its paragraph. Isolated ink above also, near numerical plus signs and a Pallas symbol, near the right-panel value 143, and above a minus sign does not establish additional notation. The interrupted closing parentheses remain parentheses. Modern fonts and reflow represent content without historical type or page identity; incidental printing and imaging intent remains uncertain.

Nachlass VI, fragment [7], retains the ascending-node sign, the low factor point between sin phi and tang one-half phi, and the fraction-height terminal point as punctuation without multiplication meaning. The calculation retains six decimal values, 54 explicit leader points, one rule and two subscript n markers. The positive cosine example +242,52 has seconds above its comma and retains its terminal point. All 146 sine terms follow the left source column before the right on each leaf. Four column starts carry full angular units and coefficient seconds marks; the other 142 terms retain sparse units. The initial delta epsilon has no added plus; the following 145 terms retain their positive signs. The first five longitude contours are Pallas. The footer signatures VII. and 70 remain. The daily secular value -0,06079 has seconds above its comma; the annual value -22,306 has none and retains its point and editorial bracket. On printed page 553, the reprint point after 0,15, the detached reprint plus tip at 5,01 and the principal tail attached to sin at 2,28 remain incidental-contour qualifications. The principal low marginal dash before the positive 0,13 term is absent from the reprint; its origin and intended status are unresolved, and it is not interpreted as a minus. On printed page 554, the reprint point before 328 and raised trace above the plus in the 0,44 term do not establish added mathematical symbols. Broken parentheses and faint specks remain qualified. Ink, type, printing and imaging remain alternatives. Modern fonts and reflow represent content rather than historical type or page identity.

Nachlass VI, fragment [8], retains the combination of the two epoch perturbations, its three written equalities and all 160 sine terms in source column order. The second sine term has Pallas without a multiplier 2 in both images and the retained leaf transcription; no multiplier is supplied. The first coefficient has no added plus, and the following 159 terms retain their positive signs. Four column starts carry coefficient seconds marks and full angular units; the remaining 156 terms retain sparse units. The integral retains upright d and italic n and t. The footer 70* remains. The daily secular value -0,06079 has seconds above its comma; the annual -22,306 has no seconds and retains its point and editorial bracket. Ten incidental contours remain qualified: the principal detached contour beside the second relation's A-prime on page 554; the reprint marginal point beside left row 19, detached plus tip at left row 39, raised trace after the final 14 at right row 13 and speck before 17,07 at right row 35 on page 555; the principal point above the gap between 234 and 3 on that page; the principal raised point between 31 and 18 at left row 9 and dot below 39 at right row 1 on page 556; and the reprint trace above the comma of 0,27 at left row 36 and point above 33 at right row 3 on that page. None establishes an additional mathematical or linguistic sign. Ink, damaged type, printing and imaging remain alternatives. Broken parentheses retain their recognized content without a claim of continuous original ink. Modern fonts and reflow represent content without historical type or page identity.

Nachlass VI, fragment [9], retains four historical paragraphs, all four differential rows and their unreduced left sides, the two epoch rows, four integrated expressions, four secular values and the complete historical footnote. The three node signs have closed lower loops. The second epoch has 121, 8 and 54,50 separated by spaces, without degree or minute marks in either image; the units supplied by the retained transcription are not inserted. Fifteen seconds-above-comma marks, six raised zeros, seven ditto signs, repeated tt, the double integral and the square on the final Pallas symbol remain. The historical note follows its original call after wiederherstellt and retains its brackets and citations to Gauss and Bessel. Six incidental contours remain qualified: the principal speck near [9], raised trace before die after 2–7, detached mark after Glieder, contour above 121 between the epoch rows, upper-right footnote-margin trace, and the reprint low trace beneath the footnote-right text. None establishes an additional linguistic or mathematical sign. Ink, damaged type, printing and imaging remain alternatives. Modern fonts and reflow represent content without historical type or page identity.

Nachlass VI, fragment [10], retains the ratio n prime to n as 7/18 in both images, the later 18/7 factor, the historical value A of 80 raised-zero 56 minutes 57 seconds, and the unmarked U row 206, 33, 46. The unit marks supplied by the retained transcription are not inserted. All three logarithmic minus-ten terms, repeated factors, twelve unbarred roots, original power positions, the complete period chain and its two M argument lists remain. Raised terminal punctuation has no multiplication meaning. Nine principal incidental contours remain qualified: the far-right speck by the first definition, the contour above Betrage, the trace after 49 seconds in the grid, the double trace above the final minus-ten term, the far-right speck beside the maximum, the trace above Das in the period paragraph, the dot below that paragraph, the dot above the first M numerator and the broad mark below the first M. None establishes an additional linguistic or mathematical sign; no index under M is supplied. Ink, damaged type, printing and imaging remain alternatives. Broken reprint parentheses are represented using the corresponding complete principal contours. Modern fonts and reflow represent content without historical type or page identity.

Nachlass VI, the first fragment [11], retains the repeated historical heading, all five tables, sparse angular unit marks and the complete bracketed note at its original call after the 1810 parenthesis. The principal scan and retained transcription support the coefficient 2,38474 in the capital-F series; the reprint's second post-comma digit is faint and broken, with a 5-like appearance. That reprint contour does not independently resolve the reading or printing intent. The node symbols, tilde-omega forms, seven ditto signs, left equation labels and all terminal punctuation retain their observed content. No unit marks are supplied in the sparse rows, including the two osculating values with only raised zeros. The retained transcription's extra stop after the final cosine term is absent from both images and is not included. Ten incidental contours remain qualified outside the historical body: the principal's far-left opening speck, slant by formula 1, low mark between F rows, short trace above the interpolation minus, far-right table-introduction mark, point beside the final node perturbation and trace above the error heading; the reprint's marks after Vergleichung and der and trace above entwickeln. None establishes an additional linguistic or mathematical sign. Ink, damaged type, printing and imaging remain alternatives. Modern fonts and reflow represent content without historical type or page identity.

Nachlass VII, letter [1] to Encke of 5 August 1816, retains the principal coefficient 10.928 with a baseline point; the reprint shows 10,928 with a comma tail. The visible contours differ, and historical printing or imaging intent remains unresolved. The point in the principal decimal is not assigned multiplication meaning. The separate point after 2 in the sine-of-45 identity retains its baseline product position. All six root signs have no overbar. The source uses psi, and the zero in x=0 has no raised degree-like zero. Both five-point omissions, all four letterspaced words, the sparse seconds-above-comma marks, the 15-degree interval and the 24 calculations are retained. Isolated marks between headings, the slant below JUPITER and detached prose specks do not establish additional linguistic signs. Modern fonts and reflow represent content and relationships without historical type or page identity.

Nachlass VII [2] retains Pallas minus Jupiter in both definitions of x; the retained leaf transcription reverses the second definition. The astronomical symbols, all prime and parenthesized indices, the two baseline product points, three right braces and three vertical rules in the argument grid retain their printed relationships. The final word Uberschrift has no visibly detached umlaut dots in the principal scan or reprint; the retained transcription has Ueberschrift. The historical note remains at its call after erste Operation. Isolated marks above the P/Q expressions, near G-prime, between bestaendig and x, and below the ending do not establish additional linguistic signs. Detached reprint strokes before h-prime and near Argument x likewise remain outside the text. Modern fonts and reflow represent content and relationships without historical type or page identity.

Nachlass VII [3] retains the staggered arguments and paired stars, the empty cells of the calculation grid, the extended rules of the Carré, and the single and double outlines of the calculation forms. The three roots are unbarred. The alpha-shaped character in the forms is represented by alpha; its a-like contour does not establish historical type identity. Partial rules retain their unequal vertical positions. Detached marginal dots, tiny strokes beside the final zeros, interruptions in border ink and the slant beside the forms do not establish additional mathematical or linguistic content. Modern fonts and reflow represent content and positional relationships, not historical type or page identity.

Nachlass VII [4] retains all fourteen tables, their sparse angular units, internal rules and grouping spaces. The shared values in A10 and A13 apply across all twenty-four rows; their braces span the complete columns. In A10, the last mark after 22 in the shared Knoten angle is represented by the single visible prime. A damaged second stroke remains a possible printing explanation. No angular-unit correction is supplied. Irregular patches, detached points and crossing strokes do not establish additional characters. Modern fonts and reflow represent the content and positional relationships without asserting historical type identity.

In Nachlass IX [1], the first table at M-prime 180, M 150 retains the qualified reading +4844. The principal scan has a 1-like final contour and the retained leaf transcription gives +4841; the reprint supplies a clearer 4. In the third table at M-prime 255, M 0, the final digit of -3379 is a qualified 9. Its open or damaged upper contour in both witnesses also permits a 6-like reading. The intended historical digit is not unambiguous. These readings do not result from numerical reconstruction.

The third table preserves the explicit +1,10644. The principal supplies the 2 in -820 at M-prime 240, M 315 and the complete +7320 at M-prime 150, M 330 where the reprint loses parts of the digits. The reprint supplies terminal digits and raised zeros clipped by the principal scan. Detached marks near entries in the three tables do not securely establish additional characters. Sparse decimal prefixes, signs, units and the different introductory and third-table numerators remain as printed. Modern fonts and reflow represent content and relative positions without asserting historical type identity.

In Nachlass IX [2], printed page 592, the constant row of the first coefficient table at M = 345 is read as 9,17751 with a lower n-marker. The principal supplies a clear first fractional digit 1 and the reprint supplies the clipped terminal digits 51 and marker. The reprint's first fractional digit is damaged and can appear 7-like at reduced scale; this is a combined-witness reading. In the sine 5M-prime row, the low sideways contour after 8,03743 at M = 285 is represented by an n rotated ninety degrees counterclockwise, and the low u-shaped contour after 7,56703 at M = 330 by a lower u-shaped glyph. These representations preserve the observed shapes and positions without deciding whether the type was rotated, inverted, a different glyph or affected by printing or imaging. They assert no additional variable or operation.

For the coefficient tables in Nachlass IX [2], the reprint supplies clipped final-column digits and raised zeros. On printed page 593, detached principal ink above the prime in the sine 5M-prime row label and after 8,72090 is absent in the reprint and is not read as additional characters. The principal supplies the legible digit in the reprint-damaged 8,98453. Modern font shapes and page reflow do not identify the original metal type.

In Nachlass IX [3], printed page 593, the principal has a detached jagged stroke between 17 and 38 in the fourth logarithmic phase. It is absent in the reprint and is not confidently identifiable as a separate character or mathematical sign. The three components 124 17 38 are retained without an added point or operation. The modern Mars and Pallas glyphs represent the observed symbols without claiming identity with the original type. The source uses unit marks in the first phase of each expression and unmarked integer triplets in the later phases; this distribution is retained.

In Nachlass IX [5], printed page 596, the numerator of the fractional exponent defining lambda is read as 3, with a qualified identification. The principal and the Springer reprint have a dense, double-lobed, 8-like contour. The faint Toronto scan has an open-left, 3-like contour with incomplete lower strokes. These witnesses support the selected exponent 3/2 without eliminating a literal 8/2, ink-filled or damaged 3, or printing and imaging alternatives. The adjacent negative fractional power and densely printed cube exponents provide comparison, but the mathematical context does not settle the source character. The historical assertion concerning the imaginary part and its inverse cosine is retained as printed. Modern font shapes and reflow do not identify the original type.

In Nachlass IX [6], printed pages 596--597, both negative fractional powers are read as minus 3/2 with a qualified identification of the numerator. On page 596, the principal and the Springer reprint have dense, 8-like contours. On page 597, the principal has a more open-left, 3-like lower contour and a dense upper lobe, while the reprint remains 8-like. This comparison supports the selected reading without resolving the page 596 contour or the related positive exponent in IX [5]. Literal 8, ink-filled or damaged 3, and printing or imaging alternatives remain possible. The mathematical context and retained transcription do not settle the original type. The coefficient-table leaders preserve the observed runs of points, including the faint fifth point after the parenthesized-zero index. Modern fonts, spacing and reflow do not establish historical type identity.

In Nachlass IX [7], printed page 600, the prime after phi in the final cosine factor is read from a descending hook joined to the upper end of the following right parenthesis in both witnesses. The same-leaf A-prime/parenthesis group and the retained transcription support the reading; a damaged or unusually drawn parenthesis without a separately intended prime remains possible. The observed contour and position are clearer than the original type or printing intent. The five root signs have no overbars. The first Limes ratio has e cos x followed by the root in its denominator, with no c base; the following cotang formula has c raised to an exponent. Both printed forms are retained. The terminal point is raised to fraction-bar height and has no multiplication meaning. Modern fonts and reflow represent the content rather than historical type identity. The positive 3/2 exponent here does not settle the related qualified readings in IX [5]--[6].

In the general introduction to the Bemerkungen on the Pallas perturbations, printed page 602, both enlarged printed witnesses read 1/1030 for the stated Jupiter mass. The retained leaf transcription instead gives 1/1050. The principal scan shows a low point after the final Manipel 2 argument, 242 28 11; that point is absent from the reprint. The point is retained without multiplication meaning, while its punctuation or printing intent remains uncertain. Other incidental ink is not treated as additional text. The date 13 December 1810 follows the clearer principal contour. Jupiter and Pallas are represented by editable modern symbol-font glyphs; those shapes and the stacked table blocks do not claim historical type or layout identity. The VII./76 printer signature belongs to the transition from printed page 601 to 602.

In the individual Bemerkungen on Pallas [I.], printed page 603, both enlarged printed witnesses show the denominator index 1 in the repeated p-double-prime ratio after “also nicht”. The retained leaf transcription omits that index. All five printed ratios, their primes and denominator indices are retained, together with the grouped numerical literal 9,791 4946 and its n marker. Modern fonts and reflow represent content rather than historical type identity.

In the individual Bemerkungen on Pallas [V.], printed page 604, the principal has a detached spot below and to the right of the final calculation-row Jupiter sign; the reprint has no corresponding spot. Its identity as a separate sign is not established. Incidental ink, printing or imaging contamination and a possible intentional low point remain alternatives; no multiplication meaning is assigned. The calculation retains all five decimal values, three seconds-above-comma forms and the rule under its third row. Its Jupiter and Pallas signs and the ascending-node sign in the final paragraph are represented with editable modern glyphs, without claiming historical type identity. The printer signature 76* belongs to the transition from printed page 603 to 604 within the joint note [III.]–[IV.].

In the opening calculation journal of Pallas Bemerkung [VI.], printed pages 604--605, the weekday symbols are retained as astronomical glyphs. The May 13 sign has the Mercury contour in both image witnesses; the retained leaf transcription has a Venus placeholder. The three-date braces, underbrace, blank cells, abbreviations and 257 individual points in 24 dot-leader runs follow the printed witnesses. Detached incidental marks near April 20, April 28, May 7--8 and the May 1--20 subtotal are not established as separate signs. Their printing or imaging origin remains uncertain. Modern glyphs, rules, braces, spacing and consecutive table blocks represent the historical content without reproducing historical type identity.

In the continuation of the calculation journal in Pallas Bemerkung [VI.], printed pages 606--607, the readings 13512, 2440, 69980, 259676, 70200 and 2204 are retained from the two image witnesses. The May 27 and November 25 signs have the Mercury contour; the retained leaf transcriptions use Venus placeholders. The blank weekday positions on May 29 and June 17 remain blank. The journal retains 124 individual points in 22 further dot-leader runs, both underbraces, the paired date and angle braces, raised zeros, brackets, abbreviations and concluding prose. Detached incidental marks near May 25, the 315 angle, 250944, the 80-day daily label and the July 13 forecast are not established as separate signs. Modern glyphs, rules, spacing and consecutive table blocks represent the historical content without reproducing historical type identity.

In the concluding paragraphs of Pallas Bemerkung [VI.], printed pages 607--609, the phase 137 75 16 and the differing sine type in the ninth term of the left series are retained. The seconds marks above decimal commas and the sparse angular units follow the source. The italic 399,07 has a comma whose lower contour is split in both image witnesses; it is represented as one comma. The historical editor's corrections and account of the manuscripts remain part of the historical text. Modern planetary glyphs and the sequential arrangement of the two series columns do not reproduce historical type identity.

In Pallas Bemerkung [IX.], printed page 609, the manuscript references Bc and Be follow their respective printed occurrences. The four angular values retain raised zeros. Brendel’s conclusion on printed pages 609--610 preserves both Latin diary quotations, with seven and five omission points, and the upright title Theoria motus. The historical credits to Encke, Westphal, Nicolai, Valentiner, Klein and Brendel remain in their source passages. Modern fonts, reflow and the closing rule represent the content rather than historical type identity or exact printing geometry.

In the Moon-theory opening, the second of the six fundamental equations retains the printed x minus x-double-prime in its final numerator. The twelve distance-denominator powers are transcribed as 3; several clogged contours resemble 8, and the identification of these powers remains qualified by the retained transcription and mathematical context. The opening Zur is upright. The terminal point after the sixth equation is raised to the fraction-bar height as punctuation, with no multiplication meaning. The part-title mark 77 with a raised asterisk remains; its function is not established here. Modern fonts and reflow represent the text rather than historical type identity.

In the conclusion of the first Moon-theory section, printed pages 615--616, the three change-of-variable identities retain both d letters in each dd t numerator. The middle identity has a qualified cubic v power; the retained leaf transcription gives a square. The first combined equation retains its printed time denominator and fourth power of p. The Greek sign in its three corresponding occurrences is psi, and tang is one operator. Cubic R and p powers and the six three-halves powers remain qualified where their contours are clogged or resemble 8. The first substituted x equation has a qualified cubic p denominator against the retained square. Its dp-squared numerator and the squared denominators under dd theta in the substitution and final equation III also have clogged contours; neighboring forms and the retained reading support 2. Ten interior low product points are distinct from the four raised terminal points, which are punctuation without multiplication meaning. The bracketed sowie and the original factor order remain. Modern glyphs, spacing, letterspacing and reflow represent content rather than historical type identity.

In the second Moon-theory section, printed pages 617--619, Greek alpha is distinct from the Latin coefficient a; repeated alpha-alpha, beta-beta and gamma-gamma products retain their form. The three forked Greek signs are psi. The resonance formulas preserve their printed signs, including the absence of a leading plus in the second expression. Its raised terminal point is punctuation without multiplication meaning, distinct from the two interior low product points. The negative three-halves powers and some small coefficient numerators retain clogged or 8-like contours. The definition of k retains the principal squared u denominator against congestion in the reprint. The final k exponent one-third is a qualified principal reading; the retained leaf gives an uncertain one-half. Context does not establish the intended printed glyph. The five logarithm relations preserve their grouping and factor order. Both identical bracketed historical notes are retained at their distinct calls. The VII./78 signature uses editable margin placement at the leaf boundary, and the final 78-star signature and closing rule remain. Modern fonts, spacing and reflow represent content rather than historical type identity.

In the opening of the third Moon-theory section, printed pages 620--622, the final psi coefficient is nine-eighths and the final solar-expansion coefficient is nine-halves, as visible in both scan witnesses. Twelve cubic powers are retained as three; several small superscripts and coefficient numerators have closed, 8-like contours, especially in the reprint. Those clogged positions remain qualified readings rather than determinations of original printing intent. Both root signs have no overbar. The three raised terminal points are punctuation without multiplication meaning. All decimal values are copied without recomputation, including the full 0,004196435 and the observational 1,0040218706. The four seconds marks are placed above their decimal commas. Both unlabelled logarithm results, the unlabelled upper-right entry and the final leading plus remain. The source wording and bracketed signed constants are retained. Detached spots are not treated as additional linguistic signs. Modern fonts, spacing and reflow represent content rather than historical type identity.

In the third Moon-theory section continuation, printed pages 622--623, both occurrences of the term two dp divided by p dv are retained. The denominator p to the fourth power is a qualified principal reading; the reprint does not show a clear exponent there, and the retained editable leaf gives three. The integration denominator retains the principal square against a three-like reprint contour. Eleven cubic powers remain qualified where small type is clogged. Two separated horizontal strokes before the three-eighths coefficient are represented by editable strokes in one operator position, without assigning them the meaning of two negations. The crossed closing delimiter order, the absent e after the second nine-fourths coefficient, all epsilon brackets and the source statement about omitting epsilon terms remain as printed. The passage ends after the integrated formula for p within a continuing historical paragraph. Modern fonts, reflow and stroke dimensions represent content rather than historical type identity.

In the third Moon-theory section ending, printed pages 624--625, the coefficients three-halves in eta-eta, p and log p, and nine-eighths in log p follow both scans rather than the differing retained editable leaf. The three-halves power of h and the final cube of h remain qualified where small type is clogged. The final u expression retains a single n in its solar term and the complete bracketed eleven-sixteenths and seventy-seven-sixteenths coefficients, whose fraction bars touch the small digits. All four numerical comparisons retain their original values, blank labels, short rules and seconds signs above decimal commas. The principal leaf ends with a rule and has no visible VII./79 signature; that signature belongs to the reprint. Modern fonts, fractions, rules and reflow represent content and positions rather than historical type identity.

In the fourth Moon-theory section opening, printed pages 626--628, both scans support nine-halves in the B-cubed/R-cubed expression and its product, and seventy-nine thirty-seconds in the unbarred-root expression; the differing retained editable readings are not adopted. Small cubic exponents and the four five-halves table numerators remain qualified where ink closes or crosses their contours. Touching fraction bars do not remove the complete numerators, denominators or factors. The two products retain low points, and the table retains its angle-column sine headings, four vertical separators, explicit zero and 79-star signature. The principal has no visible terminal mark after M in the first angle expression; a low dash-like mark in the reprint has no assigned operator meaning. Letter-spaced emphasis on Breite is retained. Modern fonts and reflow represent content and positions rather than historical type identity.

In the fourth Moon-theory section, printed pages 628--630, the numerical passage retains the printed digits, grouping, blank calculation labels, units, five rules and seconds marks. The first logarithm on page 629 has no supplied argument; the value 0,004022091 retains its zero. The following table retains thirty-two angle/coefficient pairs, repeated arguments, the single n in the nee term, four angle-column sine headings, four separators and the coefficient asterisk. Thirty-one twenty-fourths, the nine-halves and nine-eighths coefficients follow the scans. Six small cubic exponents remain qualified where ink closes their contours. A short sloping low mark after the angle 2E-lambda+2a in the principal is represented by a comma-shaped mark without mathematical operator meaning; the reprint lacks it. Its possible punctuation, incidental ink, type, printing or imaging origin remains unresolved. Detached low ink after the nee term on page 629 is not confidently a linguistic sign. Modern fonts and reflow represent content and positions rather than historical type identity.

In the fourth Moon-theory section, printed pages 630--631, the correction tables retain their bracketed headings, five and four rows, signs, repeated-letter factors and complete coefficients. The four nine-eighths entries, twenty-three twenty-fourths, sixty-five fifty-sixths and one thirty-second follow the scans. The five reductions retain both sets of Roman numerals, zeta in equation II, eta in III, all primes and the single n in the final term of II. The cubic B and R exponents remain qualified readings against closed, 8-like contours in both images; contextual corroboration does not establish the intended printed glyph. The n and theta cubes have clearer principal contours, with heavier reprint forms. Small digits and fraction-bar contacts retain their qualifications. The detached upper segment in the second correction-row closing parenthesis does not confidently establish an additional prime. Detached table ink and the principal-only tick before Folglich are not confidently linguistic signs; incidental ink, type, printing or imaging remain possible. Modern fonts and reflow represent content rather than historical type identity.

At the close of the fourth Moon-theory section, printed pages 632--633, the two pairs of dotted omission lines each contain thirty-four points per line. The comparison table retains twenty-one rows, the original coefficients and numerical values, four seconds marks above commas, sixteen Id. entries and two blank numerical positions in the first row. Its thirty-eight leader runs contain 363 points; the Mayer entry for M plus lambda retains seven points, a gap, and two points. The cubic exponents remain qualified readings where closed or heavy contours resemble eight; small three/five forms and fraction-bar contacts retain their qualifications. The reprint-only downward stroke at the first equality does not establish a new relation. Detached ink before paragraphs, between rows and below the text is not confidently linguistic. The reprint signature VII./80 is witness-specific and is not part of the principal body. Modern fonts and reflow represent content rather than historical type identity.

In the opening of the fifth Moon-theory section, printed pages 634--635, the sine and cosine tables each contain twenty-five argument/coefficient pairs, including four explicit zero coefficients across the two tables. The seven-shaped character occurs in nine positions; its encoding does not identify it as a planetary symbol. The cubic exponents of B, R and n remain qualified readings where heavy or closed contours resemble eight. Small three/five forms and fraction-bar contacts retain their qualifications. In the cosine table the coefficient for the argument 2E minus twice lambda contains three thirty-seconds, whereas the separate 4E coefficient contains thirty-three thirty-seconds. The bracketed sine argument contains plus twice M and the corresponding cosine argument minus twice M. A small speck above the reprint integration equality does not establish a different relation. Detached ink is not confidently linguistic. Modern fonts and reflow represent content rather than historical type identity.

In the fifth Moon-theory section continuation, printed pages 635--636, both coefficients twenty-seven over two hundred and fifty-six and the fourth-harmonic coefficient fifteen sixty-fourths follow the scans. The distinct constant forty-five sixty-fourths remains. The final expansion has nine-eighths and thirty-nine eighths in its two second-harmonic epsilon terms. The top addition and final expansion retain their printed closing square bracket and brace without supplying an intervening closing parenthesis. A low point after the inner v-minus-V parenthesis occurs in the principal cosine expression and is absent from the reprint; it is represented without multiplication meaning. Its printing or linguistic intent remains uncertain. The epsilon protrusion is not confidently a separate accent. Six B/R cube exponents, two three-halves powers, small three/five forms and fraction-bar contacts retain their qualifications. Six seven-shaped positions carry no asserted planetary identification. Detached ink is not treated as additional text. The 80-star signature remains. Modern fonts and reflow represent content rather than historical type identity.

In the closing tables of the fifth Moon-theory section, printed pages 637--639, the first table contains thirty-five rows, including twelve empty coefficient cells and three explicit zeros. The integrated table contains thirty-five angle/coefficient pairs and one explicit zero. Both final three-row bracket groups and the complete bracketed historical note remain. The first bracket excludes the empty third column. The opening equation has one-half in its p factor. The first table preserves nine sixty-fourths, the three-sixteenths and fifty-seven sixty-fourths terms, seven-sixteenths, Latin e beside epsilon, and two epsilons in the right-hand two-E-minus-a cell. The integrated table preserves five-sixteenths, three-sixteenths, minus seven-halves, and the opposite nine-eighths terms. The final bracket groups retain their distinct repeated-letter factors. The expression p equals h times an opening parenthesis, one and a plus sign has no supplied closing parenthesis or terminal period. Eighteen cubic exponents, including the heavily inked B/R forms, the historical-note three-halves exponent, small three/five forms and fraction-bar contacts retain their qualifications. Six seven-shaped positions carry no asserted planetary identification. Detached low ink, a separator tick and reprint specks are not confidently additional signs. Modern fonts and reflow represent content rather than historical type identity. The five printed sections end here; Gauss's historical manuscript remained unfinished, and its editorial remarks follow separately.

The historical remarks on the Moon theory consist of six paragraphs, the italic Latin diary quotation, Brendel's signature and a closing rule. They retain the two square exponents and upright differentials in du-squared over dv-squared followed by pi, repeated n and epsilon factors, complete sine arguments and the original dates and references. The negative nine thirty-seconds coefficient retains a qualified reading of its first denominator digit: the principal image has a closed, eight-like contour, the reprint is heavily inked and the retained leaf transcription reads thirty-two. The word direkte follows both printed witnesses. The historical spellings Manuskripts and Manuscript occur in their respective paragraphs. The account of the unfinished manuscript and historical editorial intervention is Brendel's. Modern fonts and reflow represent the content; they do not identify the historical type. Detached ink and reprint specks are not confidently separate signs.

The Berichtigungen und Zusätze consist of five historical correction items. They retain delta with a tilde over omega and the correction from 142 degrees 23 minutes 38 seconds to 42 degrees 23 minutes 38 seconds, including all six prime strokes and both degree signs. The corrections remain a separate historical list. F. Klein’s Bemerkungen zum siebenten Bande comprise seven paragraphs, his signature and the printer’s signature 81-star. Both leaves are unnumbered. The square brackets and large empty curly braces are part of Klein’s description of the historical edition. The original spellings and the semicolon after konnten are retained. A principal-only spot below the corrections and marks after solchen and Ausserdem in the volume note have no secure linguistic interpretation; incidental ink, type, printing and imaging remain possible explanations. Modern fonts and reflow represent the content, not historical type identity.

The historical contents retain all 158 original page references, the original headings, correspondence dates, formulae, dotted leaders and printer imprint. These references identify pages of the historical volume. They include material outside the present transcription’s stated scope. The historical dates 1811 Juni 31 and 1877 März 3 are retained as printed. The fraction beginning four-thirds and the cubic exponent in the entry for Articles 88–90 retain the qualified readings of the small printed figures. The integral retains its limits, repeated negative tt, one-half and unbarred root. The planetary symbols, original élémens and tems, and des Knoten are preserved. The principal-only low ink beneath the differential d has no secure interpretation as a separate linguistic sign. Modern fonts and reflow represent the content, not historical type identity.

The final historical publisher catalogue retains both paper editions, all fourteen prices, the printed dates, the statement that Band X and its general index were in preparation, and the full publisher, issuing institution and printer credits. Those statements describe the historical publication. The price 30 for Band VII on printing paper is a qualified reading of a figure with a closed upper loop and open lower curve. The broken figures in 25 for Band IV and 180 for Bands I--VI on writing paper retain their qualified identities. The two frames and four rules are editable; modern fonts and spacing do not reproduce historical type identity.

\Needspace{6\baselineskip}
\par\medskip
{\large Modern edition: comparison witnesses\par}\nopagebreak[4]

For the final publisher catalogue, the principal is the final unnumbered leaf of the 1906 Band VII digitization cited above. The available Springer back-matter sequence ends with the contents and does not provide this catalogue or its preceding blank leaf.



For the historical contents, the comparison witness is the Springer 1981 reprint of Band VII, \href{https://link.springer.com/content/pdf/bbm:978-3-642-92478-1/1}{back matter, printed pages 644–650}. The principal includes the Göttingen printer imprint below the final contents page; that imprint is absent from the retained reprint.



For the historical corrections and Klein’s volume note, the comparison witness is the Springer 1981 reprint of Band VII, \href{https://link.springer.com/content/pdf/bbm:978-3-642-92478-1/1}{back matter, the two unnumbered leaves before the contents}.



For the historical Moon-theory remarks, the comparison witness is the Springer 1981 reprint of Band VII, \href{https://link.springer.com/content/pdf/bbm:978-3-642-92478-1/1}{back matter, the unnumbered remarks opening and printed page 641}.



For the fifth Moon-theory section closing tables, the comparison witness is the Springer 1981 reprint of Band VII, \href{https://link.springer.com/content/pdf/bbm:978-3-642-92478-1/1}{back matter, printed pages 637--639}.



For the fifth Moon-theory section continuation, the comparison witness is the Springer 1981 reprint of Band VII, \href{https://link.springer.com/content/pdf/bbm:978-3-642-92478-1/1}{back matter, printed pages 635--636}.



For the opening of the fifth Moon-theory section, the comparison witness is the Springer 1981 reprint of Band VII, \href{https://link.springer.com/content/pdf/bbm:978-3-642-92478-1/1}{back matter, printed pages 634--635}.



For the fourth Moon-theory section ending and comparison table, the comparison witness is the Springer 1981 reprint of Band VII, \href{https://link.springer.com/content/pdf/bbm:978-3-642-92478-1/1}{back matter, printed pages 632--633}.



For the fourth Moon-theory section correction tables and five reduction equations, the comparison witness is the Springer 1981 reprint of Band VII, \href{https://link.springer.com/content/pdf/bbm:978-3-642-92478-1/1}{back matter, printed pages 630--631}.



For the fourth Moon-theory section numerical passage and continued coefficient table, the comparison witness is the Springer 1981 reprint of Band VII, \href{https://link.springer.com/content/pdf/bbm:978-3-642-92478-1/1}{back matter, printed pages 628--630}.



For the fourth Moon-theory section opening, the comparison witness is the Springer 1981 reprint of Band VII, \href{https://link.springer.com/content/pdf/bbm:978-3-642-92478-1/1}{back matter, printed pages 626--628}.



For the third Moon-theory section ending, the comparison witness is the Springer 1981 reprint of Band VII, \href{https://link.springer.com/content/pdf/bbm:978-3-642-92478-1/1}{back matter, printed pages 624--625}.



For the third Moon-theory section continuation, the comparison witness is the Springer 1981 reprint of Band VII, \href{https://link.springer.com/content/pdf/bbm:978-3-642-92478-1/1}{back matter, printed pages 622--623}.



For the third Moon-theory section opening, the comparison witness is the Springer 1981 reprint of Band VII, \href{https://link.springer.com/content/pdf/bbm:978-3-642-92478-1/1}{back matter, printed pages 620--622}.



For the second Moon-theory section, the comparison witness is the Springer 1981 reprint of Band VII, \href{https://link.springer.com/content/pdf/bbm:978-3-642-92478-1/1}{back matter, printed pages 617--619}.



For the conclusion of the first Moon-theory section, the comparison witness is the Springer 1981 reprint of Band VII, \href{https://link.springer.com/content/pdf/bbm:978-3-642-92478-1/1}{back matter, printed pages 615--616}.



For the Moon-theory part title and opening, the comparison witness is the Springer 1981 reprint of Band VII, \href{https://link.springer.com/content/pdf/bbm:978-3-642-92478-1/1}{back matter, part title, opening page and printed page 615}. Its available sequence omits the intervening blank leaf and printed page 614.



For Pallas Bemerkungen [VII.]--[IX.] and Brendel’s conclusion, the comparison witness is the Springer 1981 reprint of Band VII, \href{https://link.springer.com/content/pdf/bbm:978-3-642-92478-1/1}{back matter, printed pages 609--610}.



For the concluding paragraphs of Pallas Bemerkung [VI.], the comparison witness is the Springer 1981 reprint of Band VII, \href{https://link.springer.com/content/pdf/bbm:978-3-642-92478-1/1}{back matter, printed pages 607--609}.



For the continuation of the calculation journal in Pallas Bemerkung [VI.], the comparison witness is the Springer 1981 reprint of Band VII, \href{https://link.springer.com/content/pdf/bbm:978-3-642-92478-1/1}{back matter, printed pages 606--607}.



For the opening of Pallas Bemerkung [VI.], the comparison witness is the Springer 1981 reprint of Band VII, \href{https://link.springer.com/content/pdf/bbm:978-3-642-92478-1/1}{back matter, printed pages 604--605}.



For the individual Bemerkungen on Pallas [III.]–[V.], the comparison witness is the Springer 1981 reprint of Band VII, \href{https://link.springer.com/content/pdf/bbm:978-3-642-92478-1/1}{back matter, printed pages 603--604}.



For the individual Bemerkungen on Pallas [I.]–[II.], the comparison witness is the Springer 1981 reprint of Band VII, \href{https://link.springer.com/content/pdf/bbm:978-3-642-92478-1/1}{back matter, printed page 603}.



For the general introduction to the Bemerkungen on the Pallas perturbations, the comparison witness is the Springer 1981 reprint of Band VII, \href{https://link.springer.com/content/pdf/bbm:978-3-642-92478-1/1}{back matter, printed pages 601--603}.



For Nachlass IX [7], the equation-solving calculation and limiting ratios, the comparison witness is the Springer 1981 reprint of Band VII, \href{https://link.springer.com/content/pdf/bbm:978-3-642-92478-1/1}{back matter, printed pages 599--600}.



For Nachlass IX [6], the interpolation and coordinate series, the comparison witness is the Springer 1981 reprint of Band VII, \href{https://link.springer.com/content/pdf/bbm:978-3-642-92478-1/1}{back matter, printed pages 596--598}.



For Nachlass IX [5], the factorization and numerical example, the comparison witnesses are the Springer 1981 reprint of Band VII, \href{https://link.springer.com/content/pdf/bbm:978-3-642-92478-1/1}{back matter, printed pages 594--596}, and the University of Toronto copy, \href{https://archive.org/details/werkehrsgvonderg07gausuoft}{printed page 596}, whose faint scan is used for the fractional exponent.



For Nachlass IX [3]–[4], the perturbation expressions and the factor for the series expansions, the comparison witness is the Springer 1981 reprint of Band VII, \href{https://link.springer.com/content/pdf/bbm:978-3-642-92478-1/1}{back matter, printed pages 593--594}.



For Nachlass IX [2], the interpolation series and coefficient tables of the perturbations of Pallas by Mars, the comparison witness is the Springer 1981 reprint of Band VII, \href{https://link.springer.com/content/pdf/bbm:978-3-642-92478-1/1}{back matter, printed pages 591--593}.



For Nachlass IX [1], the introduction and three tables of the perturbations of Pallas by Mars, the comparison witness is the Springer 1981 reprint of Band VII, \href{https://link.springer.com/content/pdf/bbm:978-3-642-92478-1/1}{back matter, printed pages 587--590}.



For Nachlass VIII [5]--[6], Nicolai's letters to Gauss of 6 August 1814, 18 March and 17 July 1815, the comparison witness is the Springer 1981 reprint of Band VII, \href{https://link.springer.com/content/pdf/bbm:978-3-642-92478-1/1}{back matter, printed pages 584--586}.



For Nachlass VIII [4], Nicolai's letter to Gauss of 22 April 1814, the comparison witness is the Springer 1981 reprint of Band VII, \href{https://link.springer.com/content/pdf/bbm:978-3-642-92478-1/1}{back matter, printed pages 583--584}.



For Nachlass VIII [3], Nicolai's letter to Gauss of Good Friday 1814, the comparison witness is the Springer 1981 reprint of Band VII, \href{https://link.springer.com/content/pdf/bbm:978-3-642-92478-1/1}{back matter, printed pages 581--582}.



For Nachlass VIII [2], Nicolai's letter to Gauss of 19 March 1814, the comparison witness is the Springer 1981 reprint of Band VII, \href{https://link.springer.com/content/pdf/bbm:978-3-642-92478-1/1}{back matter, printed pages 579--581}.



For Nachlass VIII [1], the comparison witness is the Springer 1981 reprint of Band VII, \href{https://link.springer.com/content/pdf/bbm:978-3-642-92478-1/1}{back matter, printed pages 578--579}. The bracketed historical introduction belongs to the collected edition and describes Nicolai's computation. Its formulas retain the printed factors and notation. Modern fonts and reflow represent content and positional relationships without asserting historical type identity.



For Nachlass VII [4], the comparison witness is the Springer 1981 reprint of Band VII, \href{https://link.springer.com/content/pdf/bbm:978-3-642-92478-1/1}{back matter, printed pages 572--577}. The tables are compared with the 1906 principal. No independence of historical typesetting is asserted.



For both letters of Nachlass VII [3], the comparison witness is the Springer 1981 reprint of Band VII, \href{https://link.springer.com/content/pdf/bbm:978-3-642-92478-1/1}{back matter, printed pages 568--571}. The letters are compared with the 1906 principal. No independence of historical typesetting is asserted.



For Nachlass VII [2], the comparison witness is the Springer 1981 reprint of Band VII, \href{https://link.springer.com/content/pdf/bbm:978-3-642-92478-1/1}{back matter, printed pages 567--568}. The complete fragment and its historical footnote are compared with the 1906 principal. No independence of historical typesetting is asserted.



For Nachlass VII, letter [1], printed page positions 565--567, the comparison witness is the Springer 1981 reprint of Band VII, \href{https://link.springer.com/content/pdf/bbm:978-3-642-92478-1/1}{back matter, printed pages 565--567}. The principal and reprint supply the comparison for the complete letter, both trigonometric series and the root identities. No independence of historical typesetting is asserted.



For Nachlass VI, the first fragment [11], printed pages 559--561, the comparison witness is the Springer 1981 reprint of Band VII, \href{https://link.springer.com/content/pdf/bbm:978-3-642-92478-1/1}{back matter, printed pages 559--561}. Principal and reprint supply the comparison with eight oppositions, coefficient series, five tables, historical note and concluding sentence. No independence of historical typesetting is asserted.



For Nachlass VI, fragment [10], printed pages 557--559, the comparison witness is the Springer 1981 reprint of Band VII, \href{https://link.springer.com/content/pdf/bbm:978-3-642-92478-1/1}{back matter, printed pages 557--559}. Principal and reprint supply the integration, numerical grid, unmarked U row, period formula and ending. No independence of historical typesetting is asserted.



For Nachlass VI, fragment [9], printed page 557, the comparison witness is the Springer 1981 reprint of Band VII, \href{https://link.springer.com/content/pdf/bbm:978-3-642-92478-1/1}{back matter, printed page 557}. Principal and reprint supply the differential and integrated expressions, secular values, unmarked epoch units and historical note. No independence of historical typesetting is asserted.



For Nachlass VI, fragment [8], printed pages 554–556, the comparison witness is the Springer 1981 reprint of Band VII, \href{https://link.springer.com/content/pdf/bbm:978-3-642-92478-1/1}{back matter, printed pages 554–556}. Principal and reprint supply the combined epoch formula, its full sine series, sparse units and qualified contours. No independence of historical typesetting is asserted.



For Nachlass VI, fragment [7], printed pages 552–554, the comparison witness is the Springer 1981 reprint of Band VII, \href{https://link.springer.com/content/pdf/bbm:978-3-642-92478-1/1}{back matter, printed pages 552–554}. Principal and reprint are compared for the second epoch-part formula, integration calculation, full sine series, sparse units and qualified contours. No independence of historical typesetting is asserted.



For Nachlass VI, fragment [6], printed pages 550–552, the comparison witness is the Springer 1981 reprint of Band VII, \href{https://link.springer.com/content/pdf/bbm:978-3-642-92478-1/1}{back matter, printed pages 550–552}. The complete prose, integration calculation, both examples and six table panels are compared with the principal scan. No independence of historical typesetting is asserted.



For Nachlass VI, fragment [5], printed pages 549–550, the comparison witness is the Springer 1981 reprint of Band VII, \href{https://link.springer.com/content/pdf/bbm:978-3-642-92478-1/1}{back matter, printed pages 549–550}. Principal and reprint are compared for the complete prose, qualified opening fraction, integration grid and both perturbation examples. No independence of historical typesetting is asserted.



For Nachlass VI, fragment [4], printed pages 548–549, the comparison witness is the Springer 1981 reprint of Band VII, \href{https://link.springer.com/content/pdf/bbm:978-3-642-92478-1/1}{back matter, printed pages 548–549}. Principal and reprint are compared for the eccentricity formula, integration calculation, full sine series, sparse units and qualified contours. No independence of historical typesetting is asserted.



For Nachlass VI, fragment [3], printed pages 545–547, the comparison witness is the Springer 1981 reprint of Band VII, \href{https://link.springer.com/content/pdf/bbm:978-3-642-92478-1/1}{back matter, printed pages 545–547}. Principal and reprint are compared for the perihelion formulas, integration calculation, full sine series, sparse units, signatures and qualified contours. No independence of historical typesetting is asserted.



For Nachlass VI, fragment [2], printed pages 542–545, the comparison witness is the Springer 1981 reprint of Band VII, \href{https://link.springer.com/content/pdf/bbm:978-3-642-92478-1/1}{back matter, printed pages 542–545}. Principal and reprint are compared for the calculation grids, longitude signs, both sine series, sparse units and qualified contours. No independence of historical typesetting is asserted.



For Nachlass VI, fragment [1], printed pages 529–542, the comparison witness is the Springer 1981 reprint of Band VII, \href{https://link.springer.com/content/pdf/bbm:978-3-642-92478-1/1}{back matter, printed pages 529–542}. The full-width tables supply material clipped at the principal scan margins. Principal and reprint contours are also compared for damaged digits, displaced punctuation and the closing formulas. No independence of historical typesetting is asserted.



For Nachlass V, fragment [13], printed pages 526–528, the comparison witness is the Springer 1981 reprint of Band VII, \href{https://link.springer.com/content/pdf/bbm:978-3-642-92478-1/1}{back matter, printed pages 526–528}. The principal and reprint are compared for the osculating-elements table, abbreviated numerical fields, logarithmic rows, angle units and the detached contour near the fifth final phi entry. No independence of historical typesetting is asserted.



For Nachlass V, fragment [12], printed pages 525–526, the comparison witness is the Springer 1981 reprint of Band VII, \href{https://link.springer.com/content/pdf/bbm:978-3-642-92478-1/1}{back matter, printed pages 525–526}. The principal and reprint are compared for the node signs, omega accents, logarithmic block and qualified numerical contours. No independence of historical typesetting is asserted.



For Nachlass V, fragment [11], printed pages 524–525, the comparison witness is the Springer 1981 reprint of Band VII, \href{https://link.springer.com/content/pdf/bbm:978-3-642-92478-1/1}{back matter, printed pages 524–525}. It supports the pointed-loop-and-cross longitude sign, the closing parenthesis of the example and the logarithmic block. No independence of historical typesetting is asserted.



For Nachlass V, fragment [10], printed pages 522–524, the comparison witnesses include the Springer 1981 reprint of Band VII, \href{https://link.springer.com/content/pdf/bbm:978-3-642-92478-1/1}{back matter, printed pages 522–524}, and the Internet Archive scan of Band VII, camera leaf 533, printed page 523. The latter is faint and supplies limited local evidence for the two denominator readings, the logarithmic block and the contour beside 17. No independence of historical typesetting is asserted.



For Nachlass V, fragment [9], printed pages 521–522, the comparison witness is the Springer 1981 reprint of Band VII, \href{https://link.springer.com/content/pdf/bbm:978-3-642-92478-1/1}{back matter, printed pages 521–522}. It supports the irregular single-prime and double-parenthesis ending, the coefficient 37 and the logarithmic leader comparison. No independence of historical typesetting is asserted.



For Nachlass V, fragment [8], printed pages 519–521, the comparison witness is the Springer 1981 reprint of Band VII, \href{https://link.springer.com/content/pdf/bbm:978-3-642-92478-1/1}{back matter, printed pages 519–521}. It supports the comma-shaped logarithmic leader, the specified coefficient readings and the complete ending. No independence of historical typesetting is asserted.



For Nachlass V, fragment [7], printed pages 514–519, the comparison witness is the Springer 1981 reprint of Band VII, \href{https://link.springer.com/content/pdf/bbm:978-3-642-92478-1/1}{back matter, printed pages 514–519}. Its complete right-hand panel on printed page 516 supplies the part of the sine expansion clipped in the principal digitization. The remaining pages support comparisons of the coefficients, logarithmic example and final phases. No independence of historical typesetting is asserted.



For Nachlass V, fragment [6], printed pages 509–513, the comparison witness is the Springer 1981 reprint of Band VII, \href{https://link.springer.com/content/pdf/bbm:978-3-642-92478-1/1}{back matter, printed pages 509–513}. It supplies right-hand table columns clipped in the principal digitization and comparisons for qualified digits. The Internet Archive witness cited below, leaves 518–523, printed pages 508–513, supplies additional comparisons; faint, broken and gutter-obscured contours remain unresolved where stated. No independence of historical typesetting is asserted.



For Nachlass V, fragment [3], printed pages 496–500, the comparison witness is the 1906 Band VII digitization at the \href{https://lbezone.hkust.edu.hk/bib/b761099}{Hong Kong University of Science and Technology Library}, images 504–508. Its printed page 498 supplies the complete right-hand endings of roman IV, and printed page 496 supplies the comparison for the one-half coefficient. No independence of historical typesetting is asserted.

Alternative witness for Articles 98, 124 and 179, Table II and the sixth, seventh, eighth, eleventh, thirteenth and fifteenth Zusätze: \href{https://archive.org/details/werkehrsgvonderg07gausuoft}{Internet Archive, \emph{Werke}, volume 7, item werkehrsgvonderg07gausuoft}, Gerstein--University of Toronto digitization; PDF \texttt{werkehrsgvonderg07gausuoft.pdf}, leaves 140--141, printed pages 130--131, for Article 98; leaf 175, printed page 165, for Article 124. For Article 98 this witness supports the qualified exponent reading and comparison of small detached marks. For Article 124 it supplies a comparison for the qualified unprimed reading of \(X\) beside a small detached contour in the principal witness. For Article 179, leaf 255, printed page 245, supplies a faint comparison for the partially surviving exponent parentheses; it does not independently resolve their complete contours. For Table II, leaves 280 and 285, printed pages 270 and 275, support qualified readings in the historical note and two numerical entries. For the sixth Zusatz, leaf 309, printed page 299, supports the qualified opening-fraction reading and the recurrence coefficients. For the seventh Zusatz, leaves 310--311, printed pages 300--301, supply comparisons for the dense nested coefficients and final exponent; their internal contours remain partly unresolved. For the eighth Zusatz, leaf 312, printed page 302, supports the reading of the damaged closing parenthesis and double-prime in the second row of the twelve-term expression. For the eleventh Zusatz, leaf 317, printed page 307, supplies a comparison for the fragmented bracket-shaped contour after und; it does not corroborate the qualified principal reading. For the thirteenth Zusatz, leaf 319, printed page 309, supports the qualified exponent reading 3/2 in the final integral; dense principal contours remain partly unresolved. Leaves 321–322, printed pages 311–312, also support the qualified variable-glyph readings in the fifteenth Zusatz.

For the first parabolic-motion note, leaf 336, printed page 326, supplies a faint comparison for the qualified exponent reading 2 in formula [II]; it does not independently resolve that digit. For the fourth parabolic-motion note, leaf 339, printed page 329, supplies a faint comparison for the qualified first exponent reading 3/2; it does not independently resolve the dense numerator. For the sixth parabolic-motion letter, leaves 345–346, printed pages 335–336, supply comparisons for the initial fraction, root parentheses, later cosine cube, historical-note fractions and the incomplete denominator letter; the latter remains unresolved in this faint witness. For the seventh parabolic-motion note, leaves 352, 355–356 and 359, printed pages 342, 345–346 and 349, supply comparisons for the dense cubic and fractional exponents, the first series coefficient and two damaged negative indices.

For the third explanatory article of Tabula nova motus parabolici, leaf 363, printed page 353, supplies comparisons for the dense fourth powers and fractional range exponents; the latter remain partly unresolved in both witnesses. For the fifth explanatory article, leaves 364–365, printed pages 354–355, supply comparisons for six dense exponent or fractional-denominator readings of 3. These remain qualified readings supported by glyph comparison and mathematical context; the faint alternative does not resolve every contour. For the second numerical part of Tabula nova, leaf 369, printed page 359, supplies comparisons for the fragmentary final digit in argument 5,354 and the fragmentary zero in the value at 5,294; neither is independently resolved in this faint witness. Leaf 370, printed page 360, supplies comparisons for two dense zero marks, the first 9 in the value at 5,524, and the fractional digits at 5,704; the faint images do not independently resolve every contour. For the third numerical part, leaf 377, printed page 367, supplies a faint comparison for the qualified 8 at argument 32200; it does not independently resolve the incomplete upper loop. For the Encke letter, leaf 378, printed page 368, supplies a faint comparison for the two-component mark beside the first denominator A; it does not resolve its printing intent. For the letter of 1 April 1843, leaf 382, printed page 372, supplies the faint exponent comparison.

For Ceres section [2], leaf 388, printed page 378, supports the qualified cubic denominator reading and supplies a comparison for the bent last prime and fragmented closing parenthesis of the second root. The faint comparison images do not independently resolve every contour. For Article 6 of the Pallas Exposition, leaf 457, printed page 449, supplies a faint comparison of the fragmentary closing parenthesis in the first radius-sum numerator; it does not independently resolve the missing curved strokes. For Article 12 of the Pallas Exposition, leaf 466, printed page 456, supplies a faint comparison for the fragmentary parentheses in the final chain quotient; it does not independently resolve the missing contours. For Article 16 of the Pallas Exposition, leaf 470, printed page 460, supplies faint comparisons for the final denominator exponents in the force expressions; it does not resolve the qualified exponent 2 in the R expression. For Nachlass V, fragment [5], leaves 515–517, printed pages 505–507, supply comparisons for the detached pairs and qualified index readings; the faint image on page 506 does not resolve the filled nu contours. Physical independence of the typesetting or image processing is not asserted.

These editions are freely available for study and teaching.

\par\clearpage\endgroup
